% Generated by scripts/extract-theorems.py from the main book's section inputs.
% Theorem text is copied except for blue color and references to omitted material.
% External references retain their printed book numbers without links.
\documentclass[dvips]{book}

% Fonts and encoding
\usepackage[T1]{fontenc}
\usepackage{textcomp}
\renewcommand{\rmdefault}{ptm}
\usepackage[scaled=0.92]{helvet}

% Colors
\usepackage[dvipsnames, svgnames, table]{xcolor}

% Math
\usepackage{amsfonts}
\usepackage{amsmath}

% Lists, floats, figures, captions, page geometry
\usepackage{enumitem}
\usepackage[a4paper]{geometry}
\usepackage{float}
\usepackage{graphicx}
\graphicspath{{figures/}}

% Links and bookmarks
\usepackage[dvips,
  bookmarks=true,
  bookmarksdepth=3,
  colorlinks=true,
  citecolor=blue,
  urlcolor=blue,
  filecolor=blue,
  plainpages=false,
  pageanchor=true]{hyperref}

\usepackage[justification=centering, labelsep=space, skip=-10pt]{caption}
\usepackage{upref}

% Section numbering and table of contents depth
\setcounter{secnumdepth}{2}
\setcounter{tocdepth}{2}

% Book macros and environments
% ---------------------------------------------------------------------------
% Theorem-like environments
% ---------------------------------------------------------------------------
\newtheorem{theorem}{\sc\bf Theorem}
\newtheorem{example}{\sc\bf Example}
\newtheorem{definition}[theorem]{\sc\bf Definition}

\newcommand{\proof}{\noindent{\sc\bf Proof}\quad }
\newcommand{\bbox}{\phantom{1}\hfill{\rule{6pt}{6pt}}}
\newcommand{\remark}[1]{\noindent{\sc Remark}:\ #1}
\newcommand{\hint}[1]{{\sc Hint}:\ {\it #1}}
\newcommand{\note}[1]{({\sc Note}:\ #1)}

% ---------------------------------------------------------------------------
% Headings
% ---------------------------------------------------------------------------

% Unnumbered subheading, listed in the PDF bookmarks but not in the printed TOC.
\newcommand{\boxit}[1]{%
  \phantomsection
  \addcontentsline{toc}{subsubsection}{#1}%
  \subsection*{#1}}

% "Exercises" subheading at the end of a section.
\newcommand{\exercises}{%
  \phantomsection
  \pdfbookmark[2]{Exercises}{exercises.\thesection}%
  \label{exerstart:\thesection}%
  \subsection*{Exercises}%
  \noindent}

\newcommand{\technology}{\bigskip\noindent\fbox{\color{red}USING
TECHNOLOGY}\par}

% ---------------------------------------------------------------------------
% Lists
% ---------------------------------------------------------------------------
\newlist{exerciselist}{enumerate}{1}
\setlist[exerciselist]{label=\textbf{\arabic*.},ref=\arabic*}
\newlist{alist}{enumerate}{1}
\setlist[alist]{label=\textbf{(\alph*)},ref=\alph*}
\newlist{rmlist}{enumerate}{1}
\setlist[rmlist]{label=\textbf{(\roman*)},ref=\roman*}
\newlist{numberlist}{enumerate}{1}
\setlist[numberlist]{label=\textbf{\arabic*.},ref=\arabic*}

% Interrupt an exercise list with a line of text; numbering continues.
\newcommand{\exercisetext}[1]{\end{exerciselist}{\itshape #1}\begin{exerciselist}[resume]}

% ---------------------------------------------------------------------------
% Exercises and solutions
% ---------------------------------------------------------------------------
\renewcommand{\part}[1]{{\bf (#1)}}
\newcommand{\answer}[1]{\;{\bf $\mathbf{#1}$ (p.~\pageref{exer:#1})}\;}
\newcommand{\solution}{\vskip10pt\noindent{\sc\bf Solution }\  }
\newcommand{\solutionpart}[1]{\vskip10pt\noindent\underbar{\color{blue}\sc
Solution({\bf #1})\ }}

% Markers for exercises that call for computation and graphics (C, C/G)
% or for extensive use of technology (L).
\newcommand{\Cex}{\fbox{\textcolor{red}{C}}\, }
\newcommand{\CGex}{\fbox{\textcolor{red}{C/G}}\, }
\newcommand{\Lex}{\fbox{\textcolor{red}{L}}\, }

% ---------------------------------------------------------------------------
% Math shorthands
% ---------------------------------------------------------------------------
\newcommand{\dst}{\displaystyle}
\newcommand{\lims}[2]{\,\bigg|_{#1}^{#2}}

% Column vectors and matrices
\newcommand{\twocol}[2]{\left[\begin{array}{l}#1\\#2\end{array}\right]}
\newcommand{\ctwocol}[2]{\left[\begin{array}{c}#1\\#2\end{array}\right]}
\newcommand{\threecol}[3]{\left[\begin{array}{r}#1\\#2\\#3\end{array}\right]}
\newcommand{\cthreecol}[3]{\left[\begin{array}{c}#1\\#2\\#3\end{array}\right]}
\newcommand{\col}[2]{\left[\begin{array}{c}#1_1\\#1_2\\\vdots\\#1_#2\end{array}\right]}
\newcommand{\colfunc}[2]{\left[\begin{array}{c}#1_1(t)\\#1_2(t)\\\vdots\\#1_#2(t)\end{array}\right]}
\newcommand{\twobytwo}[4]{\left[\begin{array}{rr}
#1&#2\\#3&#4\end{array}\right]}
\newcommand{\threebythree}[9]{\left[\begin{array}{rrr}
#1&#2&#3\\#4&#5&#6\\#7&#8&#9\end{array}\right]}
\newcommand{\cthreebythree}[9]{\left[\begin{array}{ccc}
#1&#2&#3\\#4&#5&#6\\#7&#8&#9\end{array}\right]}
\newcommand{\matfunc}[3]{\left[\begin{array}{cccc}#1_{11}(t)&#1_{12}(t)&\cdots
&#1_{1#3}(t)\\#1_{21}(t)&#1_{22}(t)&\cdots&#1_{2#3}(t)\\\vdots&
\vdots&\ddots&\vdots\\#1_{#21}(t)&#1_{#22}(t)&\cdots&#1_{#2#3}(t)
\end{array}\right]}

% Characteristic determinant
\newcommand{\twochar}[4]{\left|\begin{array}{cc}
#1-\lambda&#2\\#3&#4-\lambda\end{array}\right|}

% ---------------------------------------------------------------------------
% Layout helpers
% ---------------------------------------------------------------------------
\newcommand{\bigbox}[1]{\vskip10pt\vbox{\hrule\hbox{\vrule\kern5pt
\vbox{\noindent
\strut\kern5pt\mbox{}\newline#1\mbox{}\newline
\kern5pt\strut}\kern5pt\vrule}\hrule}}
\newcommand{\blankbox}[1]{\vskip10pt\vbox{\hbox{\kern5pt
\vbox{\noindent\strut\kern5pt\vspace*{-15pt}#1\vspace*{-20pt}
\kern5pt\strut}\kern5pt}}}

\newenvironment{newtable}{\begin{quote}
Table\space\refstepcounter{table}\thetable.}{\end{quote}}

% Text row inside an array; must come right after a line break (\\).
\newcommand{\arraytext}[1]{\noalign{\hspace*{-1.75em}
\mbox{#1}\vspace*{-9\jot}}\nonumber\\}


% Number theorems, examples, equations, figures and tables within sections
\numberwithin{example}{section}
\numberwithin{equation}{section}
\numberwithin{theorem}{section}
\numberwithin{table}{section}
\numberwithin{figure}{section}

\hypersetup{allcolors=black}
\begin{document}

\chapter*{All Theorems}
\makeatletter
\@starttoc{toc}
\makeatother

\section*{Section 2.1}
\phantomsection
\addcontentsline{toc}{chapter}{Chapter 2}
\setcounter{chapter}{2}
\setcounter{section}{1}

\setcounter{theorem}{0}
\setcounter{equation}{11}
\begin{theorem}
  \label{thmtype:2.1.1} If $p$ is continuous on $(a,b)$, then the general solution of the
  homogeneous equation
  \begin{equation}
    \label{eq:2.1.12}y'+p(x)y=0
  \end{equation}
  on $(a,b)$ is
  \[
    y=ce^{-P(x)},
  \]
  where
  \begin{equation}
    \label{eq:2.1.13}P(x)=\int p(x)\,dx
  \end{equation}
  is any antiderivative of $p$ on $(a,b)$; that is,
  \begin{equation}
    \label{eq:2.1.14}P'(x)=p(x),\quad a<x<b.
  \end{equation}
\end{theorem}

\setcounter{theorem}{1}
\setcounter{equation}{30}
\begin{theorem}
  \label{thmtype:2.1.2} Suppose $p$ and $f$ are continuous on an open interval $(a,b)$, and let
  $y_{1}$ be any nontrivial solution of the complementary equation
  \[
    y'+p(x)y=0
  \]
  on $(a,b)$. Then:
  \begin{alist}
    \item % (a)
    The general solution of the nonhomogeneous equation
    \begin{equation}
      \label{eq:2.1.31}y'+p(x)y=f(x)
    \end{equation}
    on $(a,b)$ is
    \begin{equation}
      \label{eq:2.1.32}y=y_{1}(x)\left(c+\int f(x)/y_{1}(x)\,dx\right).
    \end{equation}

    \item % (b)
    If $x_{0}$ is an arbitrary point in $(a,b)$ and $y_{0}$ is an arbitrary real number, then the
    initial value problem
    \[
      y'+p(x)y=f(x),\quad y(x_{0})=y_{0}
    \]
    has the unique solution
    \[
      y=y_{1}(x)\left(\frac{y_{0}}{y_{1}(x_{0})}+\int^{x}_{x_{0}}\frac{f(t)}{y_{1}(t)}\,dt\right)
    \]
    on $(a,b)$.
  \end{alist}
\end{theorem}

\section*{Section 2.2}
\setcounter{chapter}{2}
\setcounter{section}{2}

\setcounter{theorem}{0}
\setcounter{equation}{9}
\begin{theorem}
  ~\label{thmtype:2.2.1} Suppose $g=g(x)$ is continuous on $(a,b)$ and $h=h(y)$ is continuous on
  $(c,d)$. Let $G$ be an antiderivative of $g$ on $(a,b)$ and let $H$ be an antiderivative of $h$ on
  $(c,d)$. Let $x_{0}$ be an arbitrary point in $(a,b)$, let $y_{0}$ be a point in $(c,d)$ such that
  $h(y_{0})\ne0$, and define
  \begin{equation}
    \label{eq:2.2.10}c=H(y_{0})-G(x_{0}).
  \end{equation}
  Then there's a function $y=y(x)$ defined on some open interval $(a_{1},b_{1})$, where
  $a\le a_{1}<x_{0}<b_{1}\le b$, such that $y(x_{0})=y_{0}$ and
  \begin{equation}
    \label{eq:2.2.11}H(y)=G(x)+c
  \end{equation}
  for $a_{1}<x<b_{1}$. Therefore $y$ is a solution of the initial value problem
  \begin{equation}
    \label{eq:2.2.12}h(y)y'=g(x),\quad y(x_{0})=y_{0}.
  \end{equation}
\end{theorem}

\section*{Section 2.3}
\setcounter{chapter}{2}
\setcounter{section}{3}

\setcounter{theorem}{0}
\setcounter{equation}{0}
\begin{theorem}
  \label{thmtype:2.3.1} \mbox{}
  \newline
  \vspace*{-5pt}
  \begin{alist}
    \item % (a)
    If $f$ is continuous on an open rectangle
    \[
      R:\{a<x<b,c<y<d\}
    \]
    that contains $(x_{0},y_{0})$ then the initial value problem
    \begin{equation}
      \label{eq:2.3.1}y'=f(x,y),\quad y(x_{0})=y_{0}
    \end{equation}
    has at least one solution on some open subinterval of $(a,b)$ that contains $x_{0}$.

    \item % (b)
    If both $f$ and $f_{y}$ are continuous on $R$ then \eqref{eq:2.3.1} has a unique solution on
    some open subinterval of $(a,b)$ that contains $x_{0}$.
  \end{alist}
\end{theorem}

\section*{Section 2.5}
\setcounter{chapter}{2}
\setcounter{section}{5}

\setcounter{theorem}{0}
\setcounter{equation}{5}
\begin{theorem}
  \label{thmtype:2.5.1} If $F=F(x,y)$ has continuous partial derivatives $F_{x}$ and $F_{y}$, then
  \begin{equation}
    \label{eq:2.5.6}F(x,y)=c\qquad\mbox{{\rm(}$c$=constant\mbox{)}},
  \end{equation}
  is an implicit solution of the differential equation
  \begin{equation}
    \label{eq:2.5.7}F_{x}(x,y)\,dx+F_{y}(x,y)\,dy=0.
  \end{equation}
\end{theorem}

\setcounter{theorem}{1}
\setcounter{equation}{11}
\begin{theorem}
  $[$The Exactness Condition$]$\label{thmtype:2.5.2} \space Suppose $M$ and $N$ are continuous and
  have continuous partial derivatives $M_{y}$ and $N_{x}$ on an open rectangle $R$. Then
  \[
    M(x,y)\,dx+N(x,y)\,dy=0
  \]
  is exact on $R$ if and only if
  \begin{equation}
    \label{eq:2.5.12}M_{y}(x,y)=N_{x}(x,y)
  \end{equation}
  for all $(x,y)$ in $R$.
\end{theorem}

\section*{Section 2.6}
\setcounter{chapter}{2}
\setcounter{section}{6}

\setcounter{theorem}{0}
\setcounter{equation}{9}
\begin{theorem}
  \label{thmtype:2.6.1} Let $M$, $N$, $M_{y}$, and $N_{x}$ be continuous on an open rectangle $R$.
  Then:
  \begin{alist}
    \item % (a)
    If $(M_{y}-N_{x})/N$ is independent of $y$ on $R$ and we define
    \[
      p(x)=\frac{M_{y}-N_{x}}{N}
    \]
    then
    \begin{equation}
      \label{eq:2.6.10}\mu(x)=\pm e^{\int p(x)\,dx}
    \end{equation}
    is an integrating factor for
    \begin{equation}
      \label{eq:2.6.11}M(x,y)\,dx+N(x,y)\,dy=0
    \end{equation}
    on $R$.

    \item % (b)
    If $(N_{x}-M_{y})/M$ is independent of $x$ on $R$ and we define
    \[
      q(y)=\frac{N_{x}-M_{y}}{M},
    \]
    then
    \begin{equation}
      \label{eq:2.6.12}\mu(y)=\pm e^{\int q(y)\,dy}
    \end{equation}
    is an integrating factor for \eqref{eq:2.6.11} on $R$.
  \end{alist}
\end{theorem}

\section*{Section 5.1}
\phantomsection
\addcontentsline{toc}{chapter}{Chapter 5}
\setcounter{chapter}{5}
\setcounter{section}{1}

\setcounter{theorem}{0}
\begin{theorem}
  \label{thmtype:5.1.1} Suppose $p$ and $q$ are continuous on an open interval $(a,b)$, let $x_{0}$
  be any point in $(a,b)$, and let $k_{0}$ and $k_{1}$ be arbitrary real numbers. Then the initial
  value problem
  \[
    y''+p(x)y'+q(x)y=0,\ y(x_{0})=k_{0},\ y'(x_{0})=k_{1}
  \]
  has a unique solution on $(a,b)$.
\end{theorem}

\setcounter{theorem}{1}
\setcounter{equation}{17}
\begin{theorem}
  \label{thmtype:5.1.2} If $y_{1}$ and $y_{2}$ are solutions of the homogeneous equation
  \begin{equation}
    \label{eq:5.1.18}y''+p(x)y'+q(x)y=0
  \end{equation}
  on $(a,b)$, then any linear combination
  \begin{equation}
    \label{eq:5.1.19}y=c_{1}y_{1}+c_{2}y_{2}
  \end{equation}
  of $y_{1}$ and $y_{2}$ is also a solution of \eqref{eq:5.1.18} on $(a,b)$.
\end{theorem}

\setcounter{theorem}{2}
\setcounter{equation}{19}
\begin{theorem}
  \label{thmtype:5.1.3} Suppose $p$ and $q$ are continuous on $(a,b)$. Then a set $\{y_{1},y_{2}\}$
  of solutions of
  \begin{equation}
    \label{eq:5.1.20}y''+p(x)y'+q(x)y=0
  \end{equation}
  on $(a,b)$ is a fundamental set if and only if $\{y_{1},y_{2}\}$ is linearly independent on
  $(a,b)$.
\end{theorem}

\setcounter{theorem}{3}
\setcounter{equation}{26}
\begin{theorem}
  \label{thmtype:5.1.4} Suppose $p$ and $q$ are continuous on $(a,b)$, let $y_{1}$ and $y_{2}$ be
  solutions of
  \begin{equation}
    \label{eq:5.1.27}y''+p(x)y'+q(x)y=0
  \end{equation}
  on $(a,b)$, and define
  \begin{equation}
    \label{eq:5.1.28}W=y_{1}y_{2}'-y_{1}'y_{2}.
  \end{equation}
  Let $x_{0}$ be any point in $(a,b)$. Then
  \begin{equation}
    \label{eq:5.1.29}W(x)=W(x_{0})e^{-\int^{x}_{x_{0}}p(t)\, dt},\quad a<x<b.
  \end{equation}
  Therefore either $W$ has no zeros in $(a,b)$ or $W\equiv0$ on $(a,b)$.
\end{theorem}

\setcounter{theorem}{4}
\setcounter{equation}{31}
\begin{theorem}
  \label{thmtype:5.1.5} Suppose $p$ and $q$ are continuous on an open interval $(a,b)$, let $y_{1}$
  and $y_{2}$ be solutions of
  \begin{equation}
    \label{eq:5.1.32}y''+p(x)y'+q(x)y=0
  \end{equation}
  on $(a,b)$, and let $W=y_{1}y_{2}'-y_{1}'y_{2}$. Then $y_{1}$ and $y_{2}$ are linearly independent
  on $(a,b)$ if and only if $W$ has no zeros on $(a,b)$.
\end{theorem}

\setcounter{theorem}{5}
\setcounter{equation}{33}
\begin{theorem}
  \label{thmtype:5.1.6} Suppose $p$ and $q$ are continuous on an open interval $(a,b)$ and let
  $y_{1}$ and $y_{2}$ be solutions of
  \begin{equation}
    \label{eq:5.1.34}y''+p(x)y'+q(x)y=0
  \end{equation}
  on $(a,b)$. Then the following statements are equivalent; that is, they are either all true or all
  false.
  \begin{alist}
    \item % (a)
    The general solution of \eqref{eq:5.1.34} on $(a,b)$ is $y=c_{1}y_{1}+c_{2}y_{2}$.

    \item % (b)
    $\{y_{1},y_{2}\}$ is a fundamental set of solutions of \eqref{eq:5.1.34} on $(a,b)$.

    \item % (c)
    $\{y_{1},y_{2}\}$ is linearly independent on $(a,b)$.

    \item % (d)
    The Wronskian of $\{y_{1},y_{2}\}$ is nonzero at some point in $(a,b)$.

    \item % (d)
    The Wronskian of $\{y_{1},y_{2}\}$ is nonzero at all points in $(a,b)$.
  \end{alist}
\end{theorem}

\setcounter{theorem}{6}
\setcounter{equation}{34}
\begin{theorem}
  \label{thmtype:5.1.7} Suppose $c$ is in $(a,b)$ and $\alpha$ and $\beta$ are real numbers, not
  both zero. Under the assumptions of Theorem~\ref{thmtype:5.1.6}, suppose $y_{1}$ and $y_{2}$ are
  solutions of \eqref{eq:5.1.34} such that
  \begin{equation}
    \label{eq:5.1.35}\alpha y_{1}(c)+\beta y_{1}'(c)=0\text{\; and\; \; }\alpha y_{2}(c)+\beta y_{2}'(c)=0.
  \end{equation}
  Then $\{y_{1},y_{2}\}$ isn't linearly independent on $(a,b)$.
\end{theorem}

\section*{Section 5.2}
\setcounter{chapter}{5}
\setcounter{section}{2}

\setcounter{theorem}{0}
\setcounter{equation}{19}
\begin{theorem}
  \label{thmtype:5.2.1} Let $a\ne0$ and let $p(r)=ar^{2}+br+c$ be the characteristic polynomial of
  \begin{equation}
    \label{eq:5.2.20}ay''+by'+cy=0.
  \end{equation}
  Then:
  \begin{alist}
    \item % (a)
    If $p(r)=0$ has distinct real roots $r_{1}$ and $r_{2}$, then the general solution of
    \eqref{eq:5.2.20} is
    \[
      y=c_{1}e^{r_{1}x}+c_{2}e^{r_{2}x}.
    \]

    \item % (b)
    If $p(r)=0$ has a repeated root $r_{1}$, then the general solution of \eqref{eq:5.2.20} is
    \[
      y=e^{r_{1}x}(c_{1}+c_{2}x).
    \]

    \item % (c)
    If $p(r)=0$ has complex conjugate roots $r_{1}=\lambda+i\omega$ and $r_{2}=\lambda-i\omega$
    (where $\omega>0$), then the general solution of \eqref{eq:5.2.20} is
    \[
      y=e^{\lambda x}(c_{1}\cos\omega x+c_{2}\sin\omega x).
    \]
  \end{alist}
\end{theorem}

\section*{Section 5.3}
\setcounter{chapter}{5}
\setcounter{section}{3}

\setcounter{theorem}{0}
\begin{theorem}
  \label{thmtype:5.3.1} Suppose $p$, $q$, and $f$ are continuous on an open interval $(a,b)$, let
  $x_{0}$ be any point in $(a,b)$, and let $k_{0}$ and $k_{1}$ be arbitrary real numbers. Then the
  initial value problem
  \[
    y''+p(x)y'+q(x)y=f(x),\quad y(x_{0})=k_{0},\quad y'(x_{0})=k_{1}
  \]
  has a unique solution on $(a,b)$.
\end{theorem}

\setcounter{theorem}{1}
\setcounter{equation}{2}
\begin{theorem}
  \label{thmtype:5.3.2} Suppose $p$, $q$, and $f$ are continuous on $(a,b)$. Let $y_{p}$ be a
  particular solution of
  \begin{equation}
    \label{eq:5.3.3}y''+p(x)y'+q(x)y=f(x)
  \end{equation}
  on $(a,b)$, and let $\{y_{1},y_{2}\}$ be a fundamental set of solutions of the complementary
  equation
  \begin{equation}
    \label{eq:5.3.4}y''+p(x)y'+q(x)y=0
  \end{equation}
  on $(a,b)$. Then $y$ is a solution of \eqref{eq:5.3.3} on $(a,b)$ if and only if
  \begin{equation}
    \label{eq:5.3.5}y=y_{p}+c_{1}y_{1}+c_{2}y_{2},
  \end{equation}
  where $c_{1}$ and $c_{2}$ are constants.
\end{theorem}

\setcounter{theorem}{2}
\begin{theorem}
  $[$The Principle of Superposition$]$ \label{thmtype:5.3.3} Suppose $y_{p_{1}}$ is a particular
  solution of
  \[
    y''+p(x)y'+q(x)y=f_{1}(x)
  \]
  on $(a,b)$ and $y_{p_{2}}$ is a particular solution of
  \[
    y''+p(x)y'+q(x)y=f_{2}(x)
  \]
  on $(a,b)$. Then
  \[
    y_{p}=y_{p_{1}}+y_{p_{2}}
  \]
  is a particular solution of
  \[
    y''+p(x)y'+q(x)y=f_{1}(x)+f_{2}(x)
  \]
  on $(a,b)$.
\end{theorem}

\section*{Section 5.5}
\setcounter{chapter}{5}
\setcounter{section}{5}

\setcounter{theorem}{0}
\setcounter{equation}{2}
\begin{theorem}
  \label{thmtype:5.5.1} Suppose $a\ne0$, $\omega$ is a positive number, and $P$ and $Q$ are polynomials, at least one nonzero. Let
  $k$ be the larger of the degrees of $P$ and $Q$. Then the equation
  \[
    ay''+by'+cy=P(x)\cos\omega x+Q(x)\sin\omega x
  \]
  has a particular solution
  \begin{equation}
    \label{eq:5.5.3}y_{p}=A(x)\cos\omega x+B(x)\sin\omega x,
  \end{equation}
  where
  \[
    A(x)=A_{0}+A_{1}x+\cdots+A_{k}x^{k}\mbox{\quad and \quad}B(x)=B_{0}+B_{1}x+\cdots+B_{k}x^{k},
  \]
  provided that $\cos\omega x$ and $\sin\omega x$ are not solutions of the complementary equation.
  The equation
  \[
    a(y''+\omega^{2}y)=P(x)\cos\omega x+Q(x)\sin\omega x
  \]
  (for which $\cos\omega x$ and $\sin\omega x$ are solutions of the complementary equation) has a particular solution
  of the form \eqref{eq:5.5.3}, where
  \[
    A(x)=A_{0}x+A_{1}x^{2}+\cdots+A_{k}x^{k+1}\mbox{\quad and \quad}B(x)=B_{0}x+B_{1}x^{2}+\cdots+B_{k}x^{k+1}.
  \]
\end{theorem}

\setcounter{theorem}{1}
\begin{theorem}
  \label{thmtype:5.5.2} If
  \[
    y_{p}=A(x)\cos\omega x+B(x)\sin\omega x,
  \]
  where $A(x)$ and $B(x)$ are polynomials with coefficients $A_{0}$ \dots, $A_{k}$ and $B_{0}$,
  \dots, $B_{k}$, then the polynomials multiplying $\sin\omega x$ in
  \[
    y_{p}',\quad y_{p}'',\quad ay_{p}''+by_{p}'+cy_{p}\mbox{\quad and \quad}y_{p}''+\omega^{2}y_{p}
  \]
  can be obtained by replacing $A_{0}$, \dots, $A_{k}$ by $B_{0}$, \dots, $B_{k}$ and $B_{0}$,
  \dots, $B_{k}$ by $-A_{0}$, \dots, $-A_{k}$ in the corresponding polynomials multiplying
  $\cos\omega x$.
\end{theorem}

\section*{Section 6.2}
\phantomsection
\addcontentsline{toc}{chapter}{Chapter 6}
\setcounter{chapter}{6}
\setcounter{section}{2}

\setcounter{theorem}{0}
\setcounter{equation}{19}
\begin{theorem}
  \label{thmtype:6.2.1} Suppose $m,k,c,F_{0}>0$ and $\omega\ge0$. Consider the amplitude $R$ of the steady state component of the
  solution of
  \[
    my''+cy'+ky=F_{0}\cos\omega t
  \]
  as a function of $\omega$.
  \newline
  \indent
  {\rm(a)} If $c\ge\sqrt{2mk}$, the maximum amplitude is $R_{\max}=F_{0}/k$ and it's attained when
  $\omega= \omega_{\max}=0$.
  \newline
  \indent
  {\rm(b)} If $c<\sqrt{2mk}$, the maximum amplitude is
  \begin{equation}
    \label{eq:6.2.20}R_{\max}=\frac{2mF_{0}}{c\sqrt{4mk-c^{2}}},
  \end{equation}
  and it's attained when
  \begin{equation}
    \label{eq:6.2.21}\omega=\omega_{\max}=\sqrt{\frac{k}{m}\left(1-\frac{c^{2}}{2km}\right)}.
  \end{equation}
\end{theorem}

\section*{Section 6.4}
\setcounter{chapter}{6}
\setcounter{section}{4}

\setcounter{theorem}{0}
\begin{theorem}
  \label{thmtype:6.4.1} The position vector of an object moving under a central force sweeps out
  equal areas in equal times. More precisely, if $t_{1}<t_{2}$, the signed area swept from
  $\theta(t_{1})$ to $\theta(t_{2})$ {\rm (Figure~6.4.3)} is given by
  \[
    A=\frac{h(t_{2}-t_{1})}{2},
  \]
  where $h=r^{2}\theta'$, which we have shown to be constant.
\end{theorem}

\section*{Section 7.1}
\phantomsection
\addcontentsline{toc}{chapter}{Chapter 7}
\setcounter{chapter}{7}
\setcounter{section}{1}

\setcounter{theorem}{1}
\begin{theorem}
  \label{thmtype:7.1.2} For any power series
  \[
    \sum_{n=0}^{\infty}a_{n}(x-x_{0})^{n},
  \]
  exactly one of these statements is true:
  \begin{rmlist}
    \item % (i)
    The power series converges only for $x=x_{0}$.

    \item % (ii)
    The power series converges for all values of $x$.

    \item % (iii)
    There's a positive number $R$ such that the power series converges if $|x-x_{0}|<R$ and diverges
    if $|x-x_{0}|>R$.
  \end{rmlist}
\end{theorem}

\setcounter{theorem}{2}
\begin{theorem}
  \label{thmtype:7.1.3} Suppose there's an integer $N$ such that $a_{n}\ne0$ if $n\ge N$ and
  \[
    \lim_{n\to\infty}\left|\frac{a_{n+1}}{a_{n}}\right|=L,
  \]
  where $0\le L\le\infty$. Then the radius of convergence of $\sum_{n=0}^{\infty}a_{n}(x-x_{0})^{n}$
  is $R=1/L$, which should be interpreted to mean that $R=0$ if $L=\infty$, or $R=\infty$ if $L=0$.
\end{theorem}

\setcounter{theorem}{3}
\setcounter{equation}{5}
\begin{theorem}
  \label{thmtype:7.1.4} A power series
  \[
    f(x)=\sum_{n=0}^{\infty}a_{n}(x-x_{0})^{n}
  \]
  with positive radius of convergence $R$ has derivatives of all orders in its open interval of
  convergence, and successive derivatives can be obtained by repeatedly differentiating term by
  term; that is,
  \begin{eqnarray}
    f'(x)      & =      & \dst{\sum_{n=1}^{\infty}na_{n}(x-x_{0})^{n-1}}\label{eq:7.1.6},      \\[10pt]
    f''(x)     & =      & \dst{\sum_{n=2}^{\infty}n(n-1)a_{n}(x-x_{0})^{n-2}},\label{eq:7.1.7} \\
               & \vdots & \nonumber                                                            \\
    f^{(k)}(x) & =      & \dst{\sum_{n=k}^{\infty}n(n-1)\cdots(n-k+1)a_{n}(x-x_{0})^{n-k}}\label{eq:7.1.8}.
  \end{eqnarray}
  Moreover, all of these series have the same radius of convergence $R$.
\end{theorem}

\setcounter{theorem}{4}
\setcounter{equation}{8}
\begin{theorem}
  \label{thmtype:7.1.5} If the power series
  \[
    f(x)=\sum_{n=0}^{\infty}a_{n}(x-x_{0})^{n}
  \]
  has a positive radius of convergence, then
  \begin{equation}
    \label{eq:7.1.9}a_{n}=\frac{f^{(n)}(x_{0})}{n!};
  \end{equation}
  that is, $\sum_{n=0}^{\infty}a_{n}(x-x_{0})^{n}$ is the Taylor series of $f$ about $x_{0}$.
\end{theorem}

\setcounter{theorem}{5}
\begin{theorem}
  \label{thmtype:7.1.6}\mbox{}
  \newline
  \begin{alist}
    \item % (a)
    If
    \[
      \sum_{n=0}^{\infty}a_{n}(x-x_{0})^{n}=\sum_{n=0}^{\infty}b_{n}(x-x_{0})^{n}
    \]
    for all $x$ in an open interval that contains $x_{0}$, then $a_{n}=b_{n}$ for $n=0$, $1$, $2$,
    \dots.

    \item % (b)
    If
    \[
      \sum_{n=0}^{\infty}a_{n}(x-x_{0})^{n}=0
    \]
    for all $x$ in an open interval that contains $x_{0}$, then $a_{n}=0$ for $n=0$, $1$, $2$,
    \dots.
  \end{alist}
\end{theorem}

\section*{Section 7.2}
\setcounter{chapter}{7}
\setcounter{section}{2}

\setcounter{theorem}{0}
\setcounter{equation}{4}
\begin{theorem}
  \label{thmtype:7.2.1} Suppose $P_{0}$, $P_{1}$, and $P_{2}$ are polynomials with no common factor
  and $P_{0}$ isn't identically zero. Let $x_{0}$ be a point such that $P_{0}(x_{0})\ne0$, and let
  $\rho$ be the distance from $x_{0}$ to the nearest zero of $P_{0}$ in the complex plane. (If
  $P_{0}$ is constant, then $\rho=\infty$.) Then every solution of
  \begin{equation}
    \label{eq:7.2.5}P_{0}(x)y''+P_{1}(x)y'+P_{2}(x)y=0
  \end{equation}
  can be represented by a power series
  \begin{equation}
    \label{eq:7.2.6}y=\sum_{n=0}^{\infty}a_{n}(x-x_{0})^{n}
  \end{equation}
  that converges at least on the open interval $(x_{0}-\rho,x_{0}+\rho)$. (If $P_{0}$ is
  nonconstant, so that $\rho$ is necessarily finite, then the open interval of convergence of
  \eqref{eq:7.2.6} may be larger than $(x_{0}-\rho,x_{0}+\rho)$. If $P_{0}$ is constant then
  $\rho=\infty$ and $(x_{0}-\rho,x_{0}+\rho)=(-\infty,\infty)$.)
\end{theorem}

\setcounter{theorem}{1}
\setcounter{equation}{22}
\begin{theorem}
  \label{thmtype:7.2.2} The coefficients $\{a_{n}\}$ in any solution
  $y=\sum_{n=0}^{\infty}a_{n}(x-x_{0})^{n}$ of
  \begin{equation}
    \label{eq:7.2.23}\left(1+\alpha(x-x_{0})^{2}\right)y''+\beta(x-x_{0})y'+\gamma y=0
  \end{equation}
  satisfy the recurrence relation
  \begin{equation}
    \label{eq:7.2.24}a_{n+2}=-\frac{p(n)}{(n+2)(n+1)}a_{n},\quad n\ge0,
  \end{equation}
  where
  \begin{equation}
    \label{eq:7.2.25}p(n)=\alpha n(n-1)+\beta n+\gamma.
  \end{equation}
  Moreover, the coefficients of the even and odd powers of $x-x_{0}$ can be computed separately as
  \begin{eqnarray}
    a_{2m+2}                                     & = & -\frac{p(2m)}{(2m+2)(2m+1)}a_{2m},\quad m\ge0\label{eq:7.2.26} \\
    \noalign{\hspace*{-15pt}\mbox{and}}\nonumber                                                                      \\
    a_{2m+3}                                     & = & -\frac{p(2m+1)}{(2m+3)(2m+2)}a_{2m+1},\quad m\ge0,\label{eq:7.2.27}
  \end{eqnarray}
  where $a_{0}$ and $a_{1}$ are arbitrary.
\end{theorem}

\section*{Section 7.4}
\setcounter{chapter}{7}
\setcounter{section}{4}

\setcounter{theorem}{2}
\setcounter{equation}{10}
\begin{theorem}
  \label{thmtype:7.4.3} Suppose the roots of the indicial equation
  \begin{equation}
    \label{eq:7.4.11}ar(r-1)+br+c=0
  \end{equation}
  are $r_{1}$ and $r_{2}$. Then the general solution of the Euler equation
  \begin{equation}
    \label{eq:7.4.12}ax^{2}y''+bxy'+cy=0
  \end{equation}
  on $(0,\infty)$ is
  \begin{eqnarray*}
    y & = & c_{1}x^{r_{1}}+c_{2}x^{r_{2}}\mbox{ if $r_{1}$ and $r_{2}$ are distinct real numbers }; \\
    y & = & x^{r_{1}}(c_{1}+c_{2}\ln x)\mbox{ if $r_{1}=r_{2}$ };                                   \\
    y & = & x^{\lambda}\left[c_{1}\cos\left(\omega\ln x\right)+c_{2}\sin\left(\omega\ln x\right)\right]\mbox{ if $r_{1},r_{2}=\lambda\pm i\omega$ with $\omega>0$}.
  \end{eqnarray*}
\end{theorem}

\section*{Section 7.5}
\setcounter{chapter}{7}
\setcounter{section}{5}

\setcounter{theorem}{0}
\setcounter{equation}{2}
\begin{theorem}
  \label{thmtype:7.5.1} Let
  \[
    Ly= x^{2}(\alpha_{0}+\alpha_{1}x+\alpha_{2}x^{2})y''+x(\beta_{0}+\beta_{1}x+\beta_{2}x^{2})y' +(\gamma_{0}+\gamma_{1}x+\gamma_{2}x^{2})y,
  \]
  and define
  \begin{eqnarray*}
    p_{0}(r)&=&\alpha_{0}r(r-1)+\beta_{0}r+\gamma_{0},\\[1\jot] p_{1}(r)&=&\alpha_{1}r(r-1)+\beta_{1}r+\gamma_{1},\\[1\jot] p_{2}(r)&=&\alpha_{2}r(r-1)+\beta_{2}r+\gamma_{2}.\\
  \end{eqnarray*}
  Suppose the series
  \begin{equation}
    \label{eq:7.5.3}y=\sum_{n=0}^{\infty}a_{n}x^{n+r}
  \end{equation}
  converges on $(0,\rho)$. Then
  \begin{equation}
    \label{eq:7.5.4}Ly=\sum_{n=0}^{\infty}b_{n}x^{n+r}
  \end{equation}
  on $(0,\rho)$, where
  \begin{equation}
    \label{eq:7.5.5}
    \begin{array}{ccl}
      b_{0} & = & p_{0}(r)a_{0},                                                       \\[1\jot]
      b_{1} & = & p_{0}(r+1)a_{1}+p_{1}(r)a_{0},                                       \\[1\jot]
      b_{n} & = & p_{0}(n+r)a_{n}+p_{1}(n+r-1)a_{n-1}+p_{2}(n+r-2)a_{n-2},\quad n\ge2.
    \end{array}
  \end{equation}
\end{theorem}

\setcounter{theorem}{1}
\setcounter{equation}{11}
\begin{theorem}
  \label{thmtype:7.5.2} Let
  \[
    Ly= x^{2}(\alpha_{0}+\alpha_{1}x+\alpha_{2}x^{2})y''+x(\beta_{0}+\beta_{1}x+\beta_{2}x^{2})y' +(\gamma_{0}+\gamma_{1}x+\gamma_{2}x^{2})y,
  \]
  where $\alpha_{0}\ne0$, and define
  \begin{eqnarray*}
    p_{0}(r)&=&\alpha_{0}r(r-1)+\beta_{0}r+\gamma_{0},\\[1\jot] p_{1}(r)&=&\alpha_{1}r(r-1)+\beta_{1}r+\gamma_{1},\\[1\jot] p_{2}(r)&=&\alpha_{2}r(r-1)+\beta_{2}r+\gamma_{2}.\\
  \end{eqnarray*}
  Suppose $r$ is a real number such that $p_{0}(n+r)$ is nonzero for all positive integers $n$.
  Define
  \begin{equation}
    \label{eq:7.5.12}
    \begin{array}{ccl}
      a_{0}(r) & = & 1,                                                                                 \\
      a_{1}(r) & = & -\dst\frac{p_{1}(r)}{p_{0}(r+1)},                                                  \\[2\jot]
      a_{n}(r) & = & -\dst\frac{p_{1}(n+r-1)a_{n-1}(r)+p_{2}(n+r-2)a_{n-2}(r)}{p_{0}(n+r)},\quad n\ge2.
    \end{array}
  \end{equation}
  Then the Frobenius series
  \begin{equation}
    \label{eq:7.5.13}y(x,r)=x^{r}\sum_{n=0}^{\infty}a_{n}(r)x^{n}
  \end{equation}
  converges and satisfies
  \begin{equation}
    \label{eq:7.5.14}Ly(x,r)=p_{0}(r)x^{r}
  \end{equation}
  on the interval $(0,\rho)$, where $\rho$ is the distance from the origin to the nearest zero of
  $A(x)=\alpha_{0}+\alpha_{1}x+\alpha_{2}x^{2}$ in the complex plane. (If $A$ is constant, then
  $\rho=\infty$.)
\end{theorem}

\setcounter{theorem}{2}
\begin{theorem}
  \label{thmtype:7.5.3} Let $L$ and $\{a_{n}(r)\}$ be as in Theorem~\ref{thmtype:7.5.2}, and
  suppose the indicial equation $p_{0}(r)=0$ of $Ly=0$ has real roots $r_{1}$ and $r_{2}$, where
  $r_{1}\ge r_{2}$. Then
  \[
    y_{1}(x)=y(x,r_{1})=x^{r_{1}}\sum_{n=0}^{\infty}a_{n}(r_{1})x^{n}
  \]
  is a Frobenius solution of $Ly=0$. Moreover, if $r_{1}-r_{2}$ isn't an integer then
  \[
    y_{2}(x)=y(x,r_{2})=x^{r_{2}}\sum_{n=0}^{\infty}a_{n}(r_{2})x^{n}
  \]
  is also a Frobenius solution of $Ly=0$, and $\{y_{1},y_{2}\}$ is a fundamental set of solutions.
\end{theorem}

\section*{Section 7.6}
\setcounter{chapter}{7}
\setcounter{section}{6}

\setcounter{theorem}{0}
\setcounter{equation}{1}
\begin{theorem}
  \label{thmtype:7.6.1} Let
  \begin{equation}
    \label{eq:7.6.2}Ly= x^{2}(\alpha_{0}+\alpha_{1}x+\alpha_{2}x^{2})y''+x(\beta_{0}+\beta_{1}x+\beta_{2}x^{2})y' +(\gamma_{0}+\gamma_{1}x+\gamma_{2}x^{2})y,
  \end{equation}
  where $\alpha_{0}\ne0$ and define
  \begin{eqnarray*}
    p_{0}(r)&=&\alpha_{0}r(r-1)+\beta_{0}r+\gamma_{0},\\[5pt] p_{1}(r)&=&\alpha_{1}r(r-1)+\beta_{1}r+\gamma_{1},\\[5pt] p_{2}(r)&=&\alpha_{2}r(r-1)+\beta_{2}r+\gamma_{2}.\\
  \end{eqnarray*}
  Suppose $r$ is a real number such that $p_{0}(n+r)$ is nonzero for all positive integers $n$, and
  define
  \[
    \begin{array}{ccl}
      a_{0}(r) & = & 1,                                                                                 \\
      a_{1}(r) & = & -\dst\frac{p_{1}(r)}{p_{0}(r+1)},                                                  \\[10pt]
      a_{n}(r) & = & -\dst\frac{p_{1}(n+r-1)a_{n-1}(r)+p_{2}(n+r-2)a_{n-2}(r)}{p_{0}(n+r)},\quad n\ge2.
    \end{array}
  \]
  Then the Frobenius series
  \begin{equation}
    \label{eq:7.6.3}y(x,r)=x^{r}\sum_{n=0}^{\infty}a_{n}(r)x^{n}
  \end{equation}
  satisfies
  \begin{equation}
    \label{eq:7.6.4}Ly(x,r)=p_{0}(r)x^{r}.
  \end{equation}
  Moreover,
  \begin{equation}
    \label{eq:7.6.5}\frac{\partial y}{\partial r}(x,r)=y(x,r)\ln x+x^{r}\sum_{n=1}^{\infty}a_{n}'(r)x^{n},
  \end{equation}
  and
  \begin{equation}
    \label{eq:7.6.6}L\left(\frac{\partial y}{\partial r}(x,r)\right)=p'_{0}(r)x^{r}+x^{r}p_{0}(r)\ln x.
  \end{equation}
\end{theorem}

\setcounter{theorem}{1}
\setcounter{equation}{7}
\begin{theorem}
  \label{thmtype:7.6.2} Let $L$ be as in Theorem~{\rm\ref{thmtype:7.6.1}} and suppose the indicial
  equation $p_{0}(r)=0$ has a repeated real root $r_{1}$. Then
  \[
    y_{1}(x)=y(x,r_{1})=x^{r_{1}}\sum_{n=0}^{\infty}a_{n}(r_{1})x^{n}
  \]
  and
  \begin{equation}
    \label{eq:7.6.8}y_{2}(x)=\frac{\partial y}{\partial r}(x,r_{1})=y_{1}(x)\ln x+x^{r_{1}}\sum_{n=1}^{\infty}a_{n}'(r_{1})x^{n}
  \end{equation}
  form a fundamental set of solutions of $Ly=0$.
\end{theorem}

\section*{Section 7.7}
\setcounter{chapter}{7}
\setcounter{section}{7}

\setcounter{theorem}{0}
\setcounter{equation}{1}
\begin{theorem}
  \label{thmtype:7.7.1} Let
  \[
    Ly= x^{2}(\alpha_{0}+\alpha_{1}x)y''+x(\beta_{0}+\beta_{1}x)y' +(\gamma_{0}+\gamma_{1}x)y,
  \]
  where $\alpha_{0}\ne0$, and define
  \begin{eqnarray*}
    p_{0}(r) & = & \alpha_{0}r(r-1)+\beta_{0}r+\gamma_{0}, \\[5pt]
    p_{1}(r) & = & \alpha_{1}r(r-1)+\beta_{1}r+\gamma_{1}.
  \end{eqnarray*}
  Suppose $r$ is a real number such that $p_{0}(n+r)$ is nonzero for all positive integers $n$, and
  define
  \begin{equation}
    \label{eq:7.7.2}
    \begin{array}{rcl}
      a_{0}(r) & = & 1,                                                          \\
      a_{n}(r) & = & -\dst\frac{p_{1}(n+r-1)}{p_{0}(n+r)}a_{n-1}(r),\quad n\ge1.
    \end{array}
  \end{equation}
  Let $r_{1}$ and $r_{2}$ be the roots of the indicial equation $p_{0}(r)=0$, and suppose
  $r_{1}=r_{2}+k$, where $k$ is a positive integer. Then
  \[
    y_{1}=x^{r_{1}}\sum_{n=0}^{\infty}a_{n}(r_{1})x^{n}
  \]
  is a Frobenius solution of $Ly=0$. Moreover, if we define
  \begin{equation}
    \label{eq:7.7.3}
    \begin{array}{rcl}
      a_{0}(r_{2}) & = & 1,                                                                              \\
      a_{n}(r_{2}) & = & -\dst\frac{p_{1}(n+r_{2}-1)}{p_{0}(n+r_{2})}a_{n-1}(r_{2}),\quad 1\le n\le k-1,
    \end{array}
  \end{equation}
  and
  \begin{equation}
    \label{eq:7.7.4}C=-\frac{p_{1}(r_{1}-1)}{k\alpha_{0}}a_{k-1}(r_{2}),
  \end{equation}
  then
  \begin{equation}
    \label{eq:7.7.5}y_{2}=x^{r_{2}}\sum_{n=0}^{k-1}a_{n}(r_{2})x^{n}+C\left(y_{1}\ln x+x^{r_{1}}\sum_{n=1}^{\infty}a_{n}'(r_{1})x^{n}\right)
  \end{equation}
  is also a solution of $Ly=0$, and $\{y_{1},y_{2}\}$ is a fundamental set of solutions.
\end{theorem}

\setcounter{theorem}{1}
\setcounter{equation}{21}
\begin{theorem}
  \label{thmtype:7.7.2} Let
  \[
    Ly= x^{2}(\alpha_{0}+\alpha_{2}x^{2})y''+x(\beta_{0}+\beta_{2}x^{2})y' +(\gamma_{0}+\gamma_{2}x^{2})y,
  \]
  where $\alpha_{0}\ne0$, and define
  \begin{eqnarray*}
    p_{0}(r) & = & \alpha_{0}r(r-1)+\beta_{0}r+\gamma_{0}, \\[5pt]
    p_{2}(r) & = & \alpha_{2}r(r-1)+\beta_{2}r+\gamma_{2}.
  \end{eqnarray*}
  Suppose $r$ is a real number such that $p_{0}(2m+r)$ is nonzero for all positive integers $m$, and
  define
  \begin{equation}
    \label{eq:7.7.22}
    \begin{array}{rcl}
      a_{0}(r)  & = & 1,                                                             \\
      a_{2m}(r) & = & -\dst\frac{p_{2}(2m+r-2)}{p_{0}(2m+r)}a_{2m-2}(r),\quad m\ge1.
    \end{array}
  \end{equation}
  Let $r_{1}$ and $r_{2}$ be the roots of the indicial equation $p_{0}(r)=0$, and suppose
  $r_{1}=r_{2}+2k$, where $k$ is a positive integer. Then
  \[
    y_{1}=x^{r_{1}}\sum_{m=0}^{\infty}a_{2m}(r_{1})x^{2m}
  \]
  is a Frobenius solution of $Ly=0$. Moreover, if we define
  \begin{eqnarray*}
    a_{0}(r_{2})  & = & 1, \\
    a_{2m}(r_{2}) & = & -\dst\frac{p_{2}(2m+r_{2}-2)}{p_{0}(2m+r_{2})}a_{2m-2}(r_{2}),\quad1\le m\le k-1
  \end{eqnarray*}
  and
  \begin{equation}
    \label{eq:7.7.23}C=-\frac{p_{2}(r_{1}-2)}{2k\alpha_{0}}a_{2k-2}(r_{2}),
  \end{equation}
  then
  \begin{equation}
    \label{eq:7.7.24}y_{2}=x^{r_{2}}\sum_{m=0}^{k-1}a_{2m}(r_{2})x^{2m}+C\left(y_{1}\ln x+x^{r_{1}}\sum_{m=1}^{\infty}a_{2m}'(r_{1})x^{2m}\right)
  \end{equation}
  is also a solution of $Ly=0$, and $\{y_{1},y_{2}\}$ is a fundamental set of solutions.
\end{theorem}

\section*{Section 8.1}
\phantomsection
\addcontentsline{toc}{chapter}{Chapter 8}
\setcounter{chapter}{8}
\setcounter{section}{1}

\setcounter{theorem}{1}
\begin{theorem}
  $[$Linearity Property$]$\label{thmtype:8.1.2} Suppose ${\cal L}(f_{i})$ is defined for $s>s_{i}$,
  $1\le i\le n$. Let $s_{0}$ be the largest of the numbers $s_{1}$, $s_{2}$, \dots,$s_{n}$, and let
  $c_{1}$, $c_{2}$,\dots, $c_{n}$ be constants. Then
  \[
    {\cal L}(c_{1}f_{1}+c_{2}f_{2}+\cdots+c_{n}f_{n})=c_{1}{\cal L}(f_{1})+c_{2}{\cal L}(f_{2}) +\cdots+c_{n}{\cal L}(f_{n})\mbox{ for }s>s_{0}.
  \]
\end{theorem}

\setcounter{theorem}{2}
\setcounter{equation}{9}
\begin{theorem}
  $[$First Shifting Theorem$]$\label{thmtype:8.1.3} If
  \begin{equation}
    \label{eq:8.1.10}F(s)=\int_{0}^{\infty}e^{-st}f(t)\,dt
  \end{equation}
  is the Laplace transform of $f(t)$ for $s>s_{0}$, then $F(s-a)$ is the Laplace transform of
  $e^{at}f(t)$ for $s>s_{0}+a$.
\end{theorem}

\setcounter{theorem}{5}
\begin{theorem}
  \label{thmtype:8.1.6} If $f$ is piecewise continuous on $[0,\infty)$ and of exponential order
  $s_{0}$, then ${\cal L}(f)$ is defined for $s>s_{0}$.
\end{theorem}

\section*{Section 8.2}
\setcounter{chapter}{8}
\setcounter{section}{2}

\setcounter{theorem}{0}
\begin{theorem}
  $[$Linearity Property$]$\label{thmtype:8.2.1} If $F_{1}$, $F_{2}$, \dots, $F_{n}$ are Laplace
  transforms and $c_{1}$, $c_{2}$, \dots, $c_{n}$ are constants, then
  \[
    {\cal L}^{-1}(c_{1}F_{1}+c_{2}F_{2}+\cdots+c_{n}F_{n})=c_{1}{\cal L}^{-1}(F_{1})+c_{2}{\cal L}^{-1}(F_{2})+\cdots+c_{n}{\cal L}^{-1}F_{n}.
  \]
\end{theorem}

\setcounter{theorem}{1}
\setcounter{equation}{5}
\begin{theorem}
  \label{thmtype:8.2.2} Suppose
  \begin{equation}
    \label{eq:8.2.6}F(s)=\frac{P(s)}{(s-s_{1})(s-s_{2})\cdots(s-s_{n})},
  \end{equation}
  where $s_{1}$, $s_{2}$, \dots, $s_{n}$ are distinct and $P$ is a polynomial of degree less than
  $n$. Then
  \[
    F(s)=\frac{A_{1}}{s-s_{1}}+\frac{A_{2}}{s-s_{2}}+\cdots+\frac{A_{n}}{s-s_{n}},
  \]
  where $A_{i}$ can be computed from \eqref{eq:8.2.6} by ignoring the factor $s-s_{i}$ and setting
  $s=s_{i}$ elsewhere.
\end{theorem}

\section*{Section 8.3}
\setcounter{chapter}{8}
\setcounter{section}{3}

\setcounter{theorem}{0}
\setcounter{equation}{0}
\begin{theorem}
  \label{thmtype:8.3.1} Suppose $f$ is continuous on $[0,\infty)$ and of exponential order $s_{0}$,
  and $f'$ is piecewise continuous on $[0,\infty)$. Then $f$ and $f'$ have Laplace transforms for
  $s>s_{0}$, and
  \begin{equation}
    \label{eq:8.3.1}{\cal L}(f')=s{\cal L}(f)-f(0).
  \end{equation}
\end{theorem}

\setcounter{theorem}{1}
\setcounter{equation}{3}
\begin{theorem}
  \label{thmtype:8.3.2} Suppose $f$ and $f'$ are continuous on $[0,\infty)$ and of exponential order
  $s_{0}$, and that $f''$ is piecewise continuous on $[0,\infty)$. Then $f$, $f'$, and $f''$ have
  Laplace transforms for $s >s_{0}$,
  \begin{equation}
    \label{eq:8.3.4}{\cal L}(f')=s{\cal L}(f)-f(0),
  \end{equation}
  and
  \begin{equation}
    \label{eq:8.3.5}{\cal L}(f'')=s^{2}{\cal L}(f)-f'(0)-sf(0).
  \end{equation}
\end{theorem}

\section*{Section 8.4}
\setcounter{chapter}{8}
\setcounter{section}{4}

\setcounter{theorem}{0}
\begin{theorem}
  \label{thmtype:8.4.1} Let $g$ be defined on $[0,\infty)$. Suppose $\tau\ge0$ and
  ${\cal L}\left(g(t+\tau)\right)$ exists for $s>s_{0}$. Then ${\cal L}\left(u(t-\tau)g(t)\right)$
  exists for $s>s_{0}$, and
  \[
    {\cal L}(u(t-\tau)g(t))=e^{-s\tau}{\cal L}\left(g(t+\tau)\right).
  \]
\end{theorem}

\setcounter{theorem}{1}
\setcounter{equation}{11}
\begin{theorem}
  $[$Second Shifting Theorem$]$\label{thmtype:8.4.2} If $\tau\ge0$ and ${\cal L}(g)$ exists for
  $s>s_{0}$ then ${\cal L}\left(u(t-\tau)g(t-\tau)\right)$ exists for $s>s_{0}$ and
  \[
    {\cal L}(u(t-\tau)g(t-\tau))=e^{-s\tau}{\cal L}(g(t)),
  \]
  or, equivalently,
  \begin{equation}
    \label{eq:8.4.12}\mbox{if }g(t)\leftrightarrow G(s),\mbox{ then }u(t-\tau)g(t-\tau)\leftrightarrow e^{-s\tau}G(s).
  \end{equation}
\end{theorem}

\section*{Section 8.5}
\setcounter{chapter}{8}
\setcounter{section}{5}

\setcounter{theorem}{0}
\begin{theorem}
  \label{thmtype:8.5.1} Suppose $a$, $b$, and $c$ are constants $(a\ne0)$, and $f$ is piecewise
  continuous on $[0,\infty)$ except for jump discontinuities at $t_{1}$, \dots, $t_{n}$, where
  \[
    0<t_{1}<\cdots<t_{n}.
  \]
  Let $k_{0}$ and $k_{1}$ be arbitrary real numbers. Then there is a unique function $y$ defined on
  $[0,\infty)$ with these properties:
  \begin{alist}
    \item % (a)
    $y(0)=k_{0}$ and $y'(0)=k_{1}$.

    \item % (b)
    $y$ and $y'$ are continuous on $[0,\infty)$.

    \item % (c)
    $y''$ is defined on every open subinterval of $[0,\infty)$ that does not contain any of the
    points $t_{1}$, \dots, $t_{n}$, and
    \[
      ay''+by'+cy=f(t)
    \]
    on every such subinterval.

    \item % (d)
    $y''$ has limits from the right and left at $t_{1}$, \dots, $t_{n}$.
  \end{alist}
\end{theorem}

\section*{Section 8.6}
\setcounter{chapter}{8}
\setcounter{section}{6}

\setcounter{theorem}{1}
\begin{theorem}
  $[$The Convolution Theorem$]$\label{thmtype:8.6.2} Suppose $f$ and $g$ are piecewise continuous on $[0,\infty)$ and their Laplace transforms converge absolutely at $s$. If ${\cal L}(f)=F$ and ${\cal L}(g)=G$, then
  \[
    {\cal L}(f*g)=FG.
  \]
\end{theorem}

\setcounter{theorem}{2}
\setcounter{equation}{12}
\begin{theorem}
  \label{thmtype:8.6.3} Suppose $f$ is continuous on $[0,\infty)$ and has a Laplace transform. Then
  the solution of the initial value problem
  \begin{equation}
    \label{eq:8.6.13}ay''+by'+cy=f(t),\quad y(0)=k_{0},\quad y'(0)=k_{1},
  \end{equation}
  is
  \begin{equation}
    \label{eq:8.6.14}y(t)=k_{0}y_{1}(t)+k_{1}y_{2}(t)+\int_{0}^{t}w(\tau)f(t-\tau)\,d\tau,
  \end{equation}
  where $y_{1}$ and $y_{2}$ satisfy
  \begin{equation}
    \label{eq:8.6.15}ay_{1}''+by_{1}'+cy_{1}=0,\quad y_{1}(0)=1,\quad y_{1}'(0)=0,
  \end{equation}
  and
  \begin{equation}
    \label{eq:8.6.16}ay_{2}''+by_{2}'+cy_{2}=0,\quad y_{2}(0)=0,\quad y_{2}'(0)=1,
  \end{equation}
  and
  \begin{equation}
    \label{eq:8.6.17}w(t)=\frac{1}{a}y_{2}(t).
  \end{equation}
\end{theorem}

\section*{Section 8.7}
\setcounter{chapter}{8}
\setcounter{section}{7}

\setcounter{theorem}{0}
\setcounter{equation}{0}
\begin{theorem}
  \label{thmtype:8.7.1} Suppose $t_{0}\ge0$. For each positive number $h$, let $y_{h}$ be the
  solution of the initial value problem
  \begin{equation}
    \label{eq:8.7.1}ay_{h}''+by_{h}'+cy_{h}=f_{h}(t),\quad y_{h}(0)=0,\quad y_{h}'(0)=0,
  \end{equation}
  where
  \begin{equation}
    \label{eq:8.7.2}f_{h}(t)=\left\{
    \begin{array}{cl}
      0,   & 0\le t<t_{0},       \\[2\jot]
      1/h, & t_{0}\le t<t_{0}+h, \\[2\jot]
      0,   & t\ge t_{0}+h,
    \end{array}\right.
  \end{equation}
  so $f_{h}$ has unit total impulse equal to the area of the shaded rectangle in
  Figure~8.7.1. Then
  \begin{equation}
    \label{eq:8.7.3}\lim_{h\to0+}y_{h}(t)=u(t-t_{0})w(t-t_{0}),
  \end{equation}
  where
  \[
    w={\cal L}^{-1}\left(\frac{1}{as^{2}+bs+c}\right).
  \]
\end{theorem}

\section*{Section 9.1}
\phantomsection
\addcontentsline{toc}{chapter}{Chapter 9}
\setcounter{chapter}{9}
\setcounter{section}{1}

\setcounter{theorem}{0}
\begin{theorem}
  \label{thmtype:9.1.1} Suppose $Ly=F$ is normal on $(a,b)$, let $x_{0}$ be a point in $(a,b)$, and
  let $k_{0}$, $k_{1}$, \dots, $k_{n-1}$ be arbitrary real numbers. Then the initial value problem
  \[
    Ly=F,\quad y(x_{0})=k_{0},\quad y'(x_{0})=k_{1},\dots,\quad y^{(n-1)}(x_{0})=k_{n-1}
  \]
  has a unique solution on $(a,b)$.
\end{theorem}

\setcounter{theorem}{1}
\begin{theorem}
  \label{thmtype:9.1.2} If $Ly=0$ is normal on $(a,b)$, then a set $\{y_{1},y_{2},\dots,y_{n}\}$ of
  $n$ solutions of $Ly=0$ on $(a,b)$ is a fundamental set if and only if it's linearly independent
  on $(a,b)$.
\end{theorem}

\setcounter{theorem}{2}
\setcounter{equation}{13}
\begin{theorem}
  \label{thmtype:9.1.3} Suppose the homogeneous linear $n$th order equation
  \begin{equation}
    \label{eq:9.1.14}P_{0}(x)y^{(n)}+P_{1}(x)y^{(n-1)}+\cdots+P_{n}(x)y=0
  \end{equation}
  is normal on $(a,b)$, let $y_{1}$, $y_{2}$, \dots, $y_{n}$ be solutions of \eqref{eq:9.1.14} on
  $(a,b)$, and let $x_{0}$ be in $(a,b)$. Then the Wronskian of $\{y_{1},y_{2},\dots,y_{n}\}$ is
  given by
  \begin{equation}
    \label{eq:9.1.15}W(x)=W(x_{0})\exp\left\{-\int^{x}_{x_{0}}\frac{P_{1}(t)}{P_{0}(t)}\, dt\right\},\quad a<x<b.
  \end{equation}
  Therefore, either $W$ has no zeros in $(a,b)$ or $W\equiv0$ on $(a,b)$.
\end{theorem}

\setcounter{theorem}{3}
\begin{theorem}
  \label{thmtype:9.1.4} Suppose $Ly=0$ is normal on $(a,b)$ and let $y_{1}$, $y_{2}$, \dots, $y_{n}$
  be $n$ solutions of $Ly=0$ on $(a,b)$. Then the following statements are equivalent; that is, they
  are either all true or all false:
  \begin{alist}
    \item % (a)
    The general solution of $Ly=0$ on $(a,b)$ is $y=c_{1}y_{1}+c_{2}y_{2}+\cdots+c_{n}y_{n}$.

    \item % (b)
    $\{y_{1},y_{2},\dots,y_{n}\}$ is a fundamental set of solutions of $Ly=0$ on $(a,b)$.

    \item % (c)
    $\{y_{1},y_{2},\dots,y_{n}\}$ is linearly independent on $(a,b)$.

    \item % (d)
    The Wronskian of $\{y_{1},y_{2},\dots,y_{n}\}$ is nonzero at some point in $(a,b)$.

    \item % (e)
    The Wronskian of $\{y_{1},y_{2},\dots,y_{n}\}$ is nonzero at all points in $(a,b)$.
  \end{alist}
\end{theorem}

\setcounter{theorem}{4}
\setcounter{equation}{15}
\begin{theorem}
  \label{thmtype:9.1.5} Suppose $Ly=0$ is normal on $(a,b)$, $c$ is in $(a,b)$, and $\alpha_{1},\alpha_{2},\dots,\alpha_{n}$ are real
  numbers, not all zero. If $y_{1},y_{2},\dots,y_{n}$ are solutions of $Ly=0$ such that
  \begin{equation}
    \label{eq:9.1.16}\alpha_{1}y_{i}(c)+\alpha_{2}y_{i}'(c)+\cdots+\alpha_{n}y_{i}^{(n-1)}(c)=0,\quad1\le i\le n,
  \end{equation}
  Then $\{y_{1},y_{2},\dots y_{n}\}$ isn't linearly independent on $(a,b)$.
\end{theorem}

\setcounter{theorem}{5}
\begin{theorem}
  \label{thmtype:9.1.6} Suppose $Ly=F$ is normal on $(a,b)$. Let $y_{p}$ be a particular solution of
  $Ly=F$ on $(a,b)$, and let $\{y_{1},y_{2},\dots,y_{n}\}$ be a fundamental set of solutions of the
  complementary equation $Ly=0$ on $(a,b)$. Then $y$ is a solution of $Ly=F$ on $(a,b)$ if and only
  if
  \[
    y=y_{p}+c_{1}y_{1}+c_{2}y_{2}+\cdots+c_{n}y_{n},
  \]
  where $c_{1},c_{2},\dots,c_{n}$ are constants.
\end{theorem}

\setcounter{theorem}{6}
\begin{theorem}
  $[$The Principle of Superposition$]$\label{thmtype:9.1.7} Suppose for each $i=1$, $2$, \dots, $r$,
  the function $y_{p_{i}}$ is a particular solution of $Ly=F_{i}$ on $(a,b)$. Then
  \[
    y_{p}=y_{p_{1}}+y_{p_{2}}+\cdots+y_{p_{r}}
  \]
  is a particular solution of
  \[
    Ly=F_{1}(x)+F_{2}(x)+\cdots+F_{r}(x)
  \]
  on $(a,b)$.
\end{theorem}

\section*{Section 9.2}
\setcounter{chapter}{9}
\setcounter{section}{2}

\setcounter{theorem}{0}
\setcounter{equation}{14}
\begin{theorem}
  \label{thmtype:9.2.1} If $m$ is a positive integer, then
  \begin{equation}
    \label{eq:9.2.15}\{e^{ax},xe^{ax},\dots,x^{m-1}e^{ax}\}
  \end{equation}
  is a fundamental set of solutions of
  \begin{equation}
    \label{eq:9.2.16}(D-a)^{m}y=0.
  \end{equation}
\end{theorem}

\setcounter{theorem}{1}
\begin{theorem}
  \label{thmtype:9.2.2} If $\omega\ne0$ and $m$ is a positive integer, then
  \[
    \begin{array}{rl}
      \{e^{\lambda x}\cos\omega x,xe^{\lambda x}\cos\omega x, & \dots,x^{m-1}e^{\lambda x}\cos\omega x,  \\
      e^{\lambda x}\sin\omega x,xe^{\lambda x}\sin\omega x,   & \dots,x^{m-1}e^{\lambda x}\sin\omega x\}
    \end{array}
  \]
  is a fundamental set of solutions of
  \[
    [(D-\lambda)^{2}+\omega^{2}]^{m}y=0.
  \]
\end{theorem}

\section*{Section 10.2}
\phantomsection
\addcontentsline{toc}{chapter}{Chapter 10}
\setcounter{chapter}{10}
\setcounter{section}{2}

\setcounter{theorem}{0}
\begin{theorem}
  \label{thmtype:10.2.1} Suppose the coefficient matrix $A$ and the forcing function ${\bf f}$ are
  continuous on $(a,b)$, let $t_{0}$ be in $(a,b)$, and let ${\bf k}$ be an arbitrary constant
  $n$-vector. Then the initial value problem
  \[
    {\bf y}'=A(t){\bf y}+{\bf f}(t),\quad{\bf y}(t_{0})={\bf k}
  \]
  has a unique solution on $(a,b)$.
\end{theorem}

\section*{Section 10.3}
\setcounter{chapter}{10}
\setcounter{section}{3}

\setcounter{theorem}{0}
\begin{theorem}
  \label{thmtype:10.3.1} Suppose the $n\times n$ matrix $A=A(t)$ is continuous on $(a,b)$. Then a
  set $\{{\bf y}_{1},{\bf y}_{2},\dots,{\bf y}_{n}\}$ of $n$ solutions of ${\bf y}'=A(t){\bf y}$ on
  $(a,b)$ is a fundamental set if and only if it's linearly independent on $(a,b)$.
\end{theorem}

\setcounter{theorem}{1}
\setcounter{equation}{5}
\begin{theorem}
  $[$Abel's Formula$]$ \label{thmtype:10.3.2} Suppose the $n\times n$ matrix $A=A(t)$ is continuous
  on $(a,b)$, let ${\bf y}_{1}$, ${\bf y}_{2}$, \dots, ${\bf y}_{n}$ be solutions of
  ${\bf y}'=A(t){\bf y}$ on $(a,b)$, and let $t_{0}$ be in $(a,b)$. Then the Wronskian of
  $\{{\bf y}_{1},{\bf y}_{2},\dots,{\bf y}_{n}\}$ is given by
  \begin{equation}
    \label{eq:10.3.6}W(t)=W(t_{0})\exp\left( \int^{t}_{t_{0}}\big[a_{11}(s)+a_{22}(s)+\cdots+a_{nn}(s)]\, ds\right),\quad a<t<b.
  \end{equation}
  Therefore, either $W$ has no zeros in $(a,b)$ or $W\equiv0$ on $(a,b)$.
\end{theorem}

\setcounter{theorem}{2}
\begin{theorem}
  \label{thmtype:10.3.3} Suppose the $n\times n$ matrix $A=A(t)$ is continuous on $(a,b)$ and let
  ${\bf y}_{1}$, ${\bf y}_{2}$, \dots,${\bf y}_{n}$ be solutions of ${\bf y}'=A(t){\bf y}$ on
  $(a,b)$. Then the following statements are equivalent; that is, they are either all true or all
  false:
  \begin{alist}
    \item % (a)
    The general solution of ${\bf y}'=A(t){\bf y}$ on $(a,b)$ is
    ${\bf y}=c_{1}{\bf y}_{1}+c_{2}{\bf y}_{2}+\cdots+c_{n}{\bf y}_{n}$, where $c_{1}$, $c_{2}$,
    \dots, $c_{n}$ are arbitrary constants.

    \item % (b)
    $\{{\bf y}_{1},{\bf y}_{2},\dots,{\bf y}_{n}\}$ is a fundamental set of solutions of
    ${\bf y}'=A(t){\bf y}$ on $(a,b)$.

    \item % (c)
    $\{{\bf y}_{1},{\bf y}_{2},\dots,{\bf y}_{n}\}$ is linearly independent on $(a,b)$.

    \item % (d)
    The Wronskian of $\{{\bf y}_{1},{\bf y}_{2},\dots,{\bf y}_{n}\}$ is nonzero at some point in
    $(a,b)$.

    \item % (e)
    The Wronskian of $\{{\bf y}_{1},{\bf y}_{2},\dots,{\bf y}_{n}\}$ is nonzero at all points in
    $(a,b)$.
  \end{alist}
\end{theorem}

\section*{Section 10.4}
\setcounter{chapter}{10}
\setcounter{section}{4}

\setcounter{theorem}{0}
\begin{theorem}
  \label{thmtype:10.4.1} Suppose the $n\times n$ constant matrix $A$ has $n$ real eigenvalues
  $\lambda_{1},\lambda_{2},\ldots,\lambda_{n}$ {\rm(}which need not be distinct{\rm)} with
  associated linearly independent eigenvectors ${\bf x}_{1},{\bf x}_{2},\ldots,{\bf x}_{n}$. Then
  the functions
  \[
    {\bf y}_{1}={\bf x}_{1}e^{\lambda_{1}t},\,{\bf y}_{2}={\bf x}_{2}e^{\lambda_{2}t},\, \dots,\,{\bf y}_{n}={\bf x}_{n}e^{\lambda_{n}t}
  \]
  form a fundamental set of solutions of ${\bf y}'=A{\bf y}$; that is, the general solution of this
  system is
  \[
    {\bf y}=c_{1}{\bf x}_{1}e^{\lambda_{1}t}+c_{2}{\bf x}_{2}e^{\lambda_{2}t}+\cdots+c_{n}{\bf x}_{n}e^{\lambda_{n}t}.
  \]
\end{theorem}

\section*{Section 10.5}
\setcounter{chapter}{10}
\setcounter{section}{5}

\setcounter{theorem}{0}
\setcounter{equation}{1}
\begin{theorem}
  \label{thmtype:10.5.1} Suppose the $n\times n$ matrix $A$ has an eigenvalue $\lambda_{1}$ of
  multiplicity $\ge2$ and the associated eigenspace has dimension $1$; that is, all
  $\lambda_{1}$-eigenvectors of $A$ are scalar multiples of an eigenvector ${\bf x}$. Then there are
  infinitely many vectors ${\bf u}$ such that
  \begin{equation}
    \label{eq:10.5.2}(A-\lambda_{1}I){\bf u}={\bf x}.
  \end{equation}
  Moreover, if ${\bf u}$ is any such vector then
  \begin{equation}
    \label{eq:10.5.3}{\bf y}_{1}={\bf x}e^{\lambda_{1}t}\quad\mbox{and }\quad{\bf y}_{2}={\bf u}e^{\lambda_{1}t}+{\bf x}te^{\lambda_{1}t}
  \end{equation}
  are linearly independent solutions of ${\bf y}'=A{\bf y}$.
\end{theorem}

\setcounter{theorem}{1}
\setcounter{equation}{8}
\begin{theorem}
  \label{thmtype:10.5.2} Suppose the $n\times n$ matrix $A$ has an eigenvalue $\lambda_{1}$ of
  multiplicity $\ge3$ and the associated eigenspace is one--dimensional; that is, all eigenvectors
  associated with $\lambda_{1}$ are scalar multiples of the eigenvector ${\bf x}$. Then there are
  infinitely many vectors ${\bf u}$ such that
  \begin{equation}
    \label{eq:10.5.9}(A-\lambda_{1}I){\bf u}={\bf x},
  \end{equation}
  and, if ${\bf u}$ is any such vector, there are infinitely many vectors ${\bf v}$ such that
  \begin{equation}
    \label{eq:10.5.10}(A-\lambda_{1}I){\bf v}={\bf u}.
  \end{equation}
  If ${\bf u}$ satisfies {\rm\eqref{eq:10.5.9}} and ${\bf v}$ satisfies {\rm\eqref{eq:10.5.10}},
  then
  \begin{eqnarray*}
    {\bf y}_{1} & = & {\bf x}e^{\lambda_{1}t},                                      \\
    {\bf y}_{2} & = & {\bf u}e^{\lambda_{1}t}+{\bf x}te^{\lambda_{1}t},\mbox{ and } \\
    {\bf y}_{3} & = & {\bf v}e^{\lambda_{1}t}+{\bf u}te^{\lambda_{1}t}+{\bf x}{\frac{t^{2}e^{\lambda_{1}t}}{2}}
  \end{eqnarray*}
  are linearly independent solutions of ${\bf y}'=A{\bf y}$.
\end{theorem}

\setcounter{theorem}{2}
\setcounter{equation}{11}
\begin{theorem}
  \label{thmtype:10.5.3} Suppose the $n\times n$ matrix $A$ has an eigenvalue $\lambda_{1}$ of
  multiplicity $\ge3$ and the associated eigenspace is two--dimensional; that is, all eigenvectors
  of $A$ associated with $\lambda_{1}$ are linear combinations of two linearly independent
  eigenvectors ${\bf x}_{1}$ and ${\bf x}_{2}$. Then there are constants $\alpha$ and $\beta$ (not
  both zero) such that if
  \begin{equation}
    \label{eq:10.5.12}{\bf x}_{3}=\alpha{\bf x}_{1}+\beta{\bf x}_{2},
  \end{equation}
  then there are infinitely many vectors ${\bf u}$ such that
  \begin{equation}
    \label{eq:10.5.13}(A-\lambda_{1}I){\bf u}={\bf x}_{3}.
  \end{equation}
  If ${\bf u}$ satisfies {\rm\eqref{eq:10.5.13}}, then
  \begin{eqnarray}
    {\bf y}_{1} & = & {\bf x}_{1}e^{\lambda_{1}t},\nonumber            \\
    {\bf y}_{2} & = & {\bf x}_{2}e^{\lambda_{1}t},\mbox{and }\nonumber \\
    {\bf y}_{3} & = & {\bf u}e^{\lambda_{1}t}+{\bf x}_{3}te^{\lambda_{1}t}\label{eq:10.5.14},
  \end{eqnarray}
  are linearly independent solutions of ${\bf y}'=A{\bf y}$.
\end{theorem}

\section*{Section 10.6}
\setcounter{chapter}{10}
\setcounter{section}{6}

\setcounter{theorem}{0}
\setcounter{equation}{2}
\begin{theorem}
  \label{thmtype:10.6.1} Let $A$ be an $n\times n$ matrix with real entries. Let
  $\lambda=\alpha+i\beta$ {\rm(}$\beta\ne0${\rm}) be a complex eigenvalue of $A$ and let
  ${\bf x}={\bf u}+i{\bf v}$ be an associated eigenvector, where ${\bf u}$ and ${\bf v}$ have real
  components. Then ${\bf u}$ and ${\bf v}$ are both nonzero and
  \[
    {\bf y}_{1}=e^{\alpha t}({\bf u}\cos\beta t-{\bf v}\sin\beta t) \mbox{\quad and\quad}{\bf y}_{2}=e^{\alpha t}({\bf u}\sin\beta t+{\bf v}\cos\beta t),
  \]
  which are the real and imaginary parts of
  \begin{equation}
    \label{eq:10.6.3}e^{\alpha t}(\cos\beta t+i\sin\beta t)({\bf u}+i{\bf v}),
  \end{equation}
  are linearly independent solutions of ${\bf y}'=A{\bf y}$.
\end{theorem}

\section*{Section 10.7}
\setcounter{chapter}{10}
\setcounter{section}{7}

\setcounter{theorem}{0}
\begin{theorem}
  \label{thmtype:10.7.1} Suppose the $n\times n$ matrix function $A$ and the $n$-vector function
  ${\bf f}$ are continuous on $(a,b)$. Let ${\bf y}_{p}$ be a particular solution of
  ${\bf y}'=A(t){\bf y}+{\bf f}(t)$ on $(a,b)$, and let
  $\{{\bf y}_{1},{\bf y}_{2},\dots,{\bf y}_{n}\}$ be a fundamental set of solutions of the
  complementary equation ${\bf y}'=A(t){\bf y}$ on $(a,b)$. Then ${\bf y}$ is a solution of
  ${\bf y}'=A(t){\bf y}+{\bf f}(t)$ on $(a,b)$ if and only if
  \[
    {\bf y}={\bf y}_{p}+c_{1}{\bf y}_{1}+c_{2}{\bf y}_{2}+\cdots+c_{n}{\bf y}_{n},
  \]
  where $c_{1}$, $c_{2}$, \dots, $c_{n}$ are constants.
\end{theorem}

\section*{Section 11.1}
\phantomsection
\addcontentsline{toc}{chapter}{Chapter 11}
\setcounter{chapter}{11}
\setcounter{section}{1}

\setcounter{theorem}{0}
\begin{theorem}
  \label{thmtype:11.1.1} Problems $1$--$5$ have no negative eigenvalues. Moreover, $\lambda=0$ is an
  eigenvalue of Problems $2$ and $5$, with associated eigenfunction $y_{0}=1$, but $\lambda=0$ isn't
  an eigenvalue of Problems~$1$, $3$, or $4$.
\end{theorem}

\setcounter{theorem}{1}
\begin{theorem}
  \label{thmtype:11.1.2} The eigenvalue problem
  \[
    y''+\lambda y=0,\quad y(0)=0,\quad y(L)=0
  \]
  has infinitely many positive eigenvalues $\lambda_{n}=n^{2}\pi^{2}/L^{2}$, with associated
  eigenfunctions
  \[
    y_{n}=\sin\frac{n\pi x}{L},\quad n=1,2,3,\dots.
  \]
  There are no other eigenvalues.
\end{theorem}

\setcounter{theorem}{2}
\begin{theorem}
  \label{thmtype:11.1.3} The eigenvalue problem
  \[
    y''+\lambda y=0,\quad y'(0)=0,\quad y'(L)=0
  \]
  has the eigenvalue $\lambda_{0}=0$, with associated eigenfunction $y_{0}=1$, and infinitely many
  positive eigenvalues $\lambda_{n}=n^{2}\pi^{2}/L^{2}$, with associated eigenfunctions
  \[
    y_{n}=\cos\frac{n\pi x}{L},n=1,2,3\dots.
  \]
  There are no other eigenvalues.
\end{theorem}

\setcounter{theorem}{3}
\begin{theorem}
  \label{thmtype:11.1.4} The eigenvalue problem
  \[
    y''+\lambda y=0,\quad y(0)=0,\quad y'(L)=0
  \]
  has infinitely many positive eigenvalues $\lambda_{n}=(2n-1)^{2}\pi^{2}/4L^{2}$, with associated
  eigenfunctions
  \[
    y_{n}=\sin\frac{(2n-1)\pi x}{2L},\quad n=1,2,3,\dots.
  \]
  There are no other eigenvalues.
\end{theorem}

\setcounter{theorem}{4}
\begin{theorem}
  \label{thmtype:11.1.5} The eigenvalue problem
  \[
    y''+\lambda y=0,\quad y'(0)=0,\quad y(L)=0
  \]
  has infinitely many positive eigenvalues $\lambda_{n}=(2n-1)^{2}\pi^{2}/4L^{2}$, with associated
  eigenfunctions
  \[
    y_{n}=\cos\frac{(2n-1)\pi x}{2L},\quad n=1,2,3,\dots.
  \]
  There are no other eigenvalues.
\end{theorem}

\setcounter{theorem}{5}
\begin{theorem}
  \label{thmtype:11.1.6} The eigenvalue problem
  \[
    y''+\lambda y=0,\quad y(-L)=y(L),\quad y'(-L)=y'(L),
  \]
  has the eigenvalue $\lambda_{0}=0$, with associated eigenfunction $y_{0}=1$ and infinitely many
  positive eigenvalues $\lambda_{n}=n^{2}\pi^{2}/L^{2}$, with associated eigenfunctions
  \[
    y_{1n}=\cos\frac{n\pi x}{L}\mbox{\quad and \quad}y_{2n}=\sin\frac{n\pi x}{L},\quad n=1,2,3,\dots.
  \]
  There are no other eigenvalues.
\end{theorem}

\section*{Section 11.2}
\setcounter{chapter}{11}
\setcounter{section}{2}

\setcounter{theorem}{0}
\setcounter{equation}{0}
\begin{theorem}
  \label{thmtype:11.2.1} Suppose the functions $\phi_{1}$, $\phi_{2}$, $\phi_{3}$, \dots, are
  orthogonal on $[a,b]$ and
  \begin{equation}
    \label{eq:11.2.1}\int_{a}^{b}\phi_{n}^{2}(x)\,dx\ne0,\quad n=1,2,3,\dots.
  \end{equation}
  Let $c_{1}$, $c_{2}$, $c_{3}$,\dots be constants such that the partial sums
  $f_{N}(x)=\sum_{m=1}^{N}c_{m}\phi_{m}(x)$ satisfy the inequalities
  \[
    |f_{N}(x)|\le M,\quad a\le x\le b,\quad N=1,2,3,\dots
  \]
  for some constant $M<\infty$. Suppose also that the series
  \begin{equation}
    \label{eq:11.2.2}f(x)=\sum_{m=1}^{\infty}c_{m}\phi_{m}(x)
  \end{equation}
  converges and is integrable on $[a,b]$. Then
  \begin{equation}
    \label{eq:11.2.3}c_{n}=\frac{\dst\int_{a}^{b}f(x)\phi_{n}(x)\,dx}{\dst\int_{a}^{b}\phi_{n}^{2}(x)\,dx},\quad n=1,2,3,\dots.
  \end{equation}
\end{theorem}

\setcounter{theorem}{3}
\setcounter{equation}{7}
\begin{theorem}
  \label{thmtype:11.2.4} If $f$ is piecewise smooth on $[-L,L]$, then the Fourier series
  \begin{equation}
    \label{eq:11.2.8}F(x)=a_{0}+\sum_{n=1}^{\infty}\left(a_{n}\cos\frac{n\pi x}{L}+b_{n}\sin\frac{n\pi x}{L}\right)
  \end{equation}
  of $f$ on $[-L,L]$ converges for all $x$ in $[-L,L]$; moreover,
  \[
    F(x)= \left\{
    \begin{array}{cl}
      f(x)                       & \mbox{if $-L<x<L$ and $f$ is continuous at $x$ }    \\[2\jot]
      \dst\frac{f(x-)+f(x+)}{2}  & \mbox{if $-L<x<L$ and $f$ is discontinuous at $x$ } \\[2\jot]
      \dst\frac{f(-L+)+f(L-)}{2} & \mbox{if $x=L$ or $x=-L$.}
    \end{array}\right.
  \]
\end{theorem}

\setcounter{theorem}{4}
\begin{theorem}
  \label{thmtype:11.2.5} Suppose $u$ is even and $v$ is odd on $[-L,L]$. Then:
  \[
    \mbox{\part{a}\quad}\int_{-L}^{L}u(x)\,dx=2\int_{0}^{L}u(x)\,dx, \mbox{\quad \part{b}\quad}\int_{-L}^{L}v(x)\,dx=0,
  \]
  \[
    \mbox{\part{c}\quad}\int_{-L}^{L}u(x)\cos\frac{n\pi x}{L}\,dx=2\int_{0}^{L}u(x)\cos\frac{n\pi x}{L}\,dx,
  \]
  \[
    \mbox{\part{d}\quad}\int_{-L}^{L}v(x)\sin\frac{n\pi x}{L}\,dx= 2\int_{0}^{L}v(x)\sin\frac{n\pi x}{L}\,dx,
  \]
  \[
    \mbox{\part{e}\quad}\int_{-L}^{L}u(x)\sin\frac{n\pi x}{L}\,dx=0 \mbox{\quad and \quad \part{f}\quad}\int_{-L}^{L}v(x)\cos\frac{n\pi x}{L}\,dx=0.
  \]
\end{theorem}

\setcounter{theorem}{5}
\begin{theorem}
  \label{thmtype:11.2.6} Suppose $f$ is integrable on $[-L,L]$.
  \begin{alist}
    \item % (a)
    If $f$ is even, the Fourier series of $f$ on $[-L,L]$ is
    \[
      F(x)=a_{0}+\sum_{n=1}^{\infty}a_{n}\cos\frac{n\pi x}{L},
    \]
    where
    \[
      a_{0}=\frac{1}{L}\int_{0}^{L}f(x)\,dx\mbox{\quad and \quad}a_{n}=\frac{2}{L}\int_{0}^{L}f(x)\cos\frac{n\pi x}{L}\,dx,\quad n\ge1.
    \]

    \item % (b)
    If $f$ is odd, the Fourier series of $f$ on $[-L,L]$ is
    \[
      F(x)=\sum_{n=1}^{\infty}b_{n}\sin\frac{n\pi x}{L},
    \]
    where
    \[
      b_{n}=\frac{2}{L}\int_{0}^{L}f(x)\sin\frac{n\pi x}{L}\,dx.
    \]
  \end{alist}
\end{theorem}

\section*{Section 11.3}
\setcounter{chapter}{11}
\setcounter{section}{3}

\setcounter{theorem}{0}
\begin{theorem}
  \label{thmtype:11.3.1} If $f$ is piecewise smooth on $[0,L]$, then the Fourier cosine series
  \[
    C(x)=a_{0}+\sum_{n=1}^{\infty}a_{n}\cos\frac{n\pi x}{L}
  \]
  of $f$ on $[0,L]$, with
  \[
    a_{0}=\frac{1}{L}\int_{0}^{L}f(x)\,dx \mbox{\quad and \quad}a_{n}=\frac{2}{L}\int_{0}^{L}f(x)\cos\frac{n\pi x}{L}\,dx,\quad n=1,2,3,\dots,
  \]
  converges for all $x$ in $[0,L]$; moreover,
  \[
    C(x)= \left\{
    \begin{array}{cl}
      f(0+)                     & \mbox{if }x=0                                     \\
      f(x)                      & \mbox{if $0<x<L$ and $f$ is continuous at $x$}    \\[2\jot]
      \dst\frac{f(x-)+f(x+)}{2} & \mbox{if $0<x<L$ and $f$ is discontinuous at $x$} \\[2\jot]
      f(L-)                     & \mbox{if }x=L.
    \end{array}\right.
  \]
\end{theorem}

\setcounter{theorem}{1}
\begin{theorem}
  \label{thmtype:11.3.2} If $f$ is piecewise smooth on $[0,L]$, then the Fourier sine series
  \[
    S(x)=\sum_{n=1}^{\infty}b_{n}\sin\frac{n\pi x}{L}
  \]
  of $f$ on $[0,L]$, with
  \[
    b_{n}=\frac{2}{L}\int_{0}^{L}f(x)\sin\frac{n\pi x}{L}\,dx,
  \]
  converges for all $x$ in $[0,L]$; moreover,
  \[
    S(x)= \left\{
    \begin{array}{cl}
      0                         & \mbox{if }x=0                                     \\
      f(x)                      & \mbox{if $0<x<L$ and $f$ is continuous at $x$}    \\[2\jot]
      \dst\frac{f(x-)+f(x+)}{2} & \mbox{if $0<x<L$ and $f$ is discontinuous at $x$} \\[2\jot]
      0                         & \mbox{if }x=L.
    \end{array}\right.
  \]
\end{theorem}

\setcounter{theorem}{2}
\begin{theorem}
  \label{thmtype:11.3.3} If $f$ is piecewise smooth on $[0,L]$, then the mixed Fourier cosine series
  \[
    C_{M}(x)=\sum_{n=1}^{\infty}c_{n}\cos\frac{(2n-1)\pi x}{2L}
  \]
  of $f$ on $[0,L]$, with
  \[
    c_{n}=\frac{2}{L}\int_{0}^{L}f(x)\cos\frac{(2n-1)\pi x}{2L}\,dx,
  \]
  converges for all $x$ in $[0,L]$; moreover,
  \[
    C_{M}(x)= \left\{
    \begin{array}{cl}
      f(0+)                     & \mbox{if $x=0$}                                    \\
      f(x)                      & \mbox{if $0<x<L$ and $f$ is continuous at $x$ }    \\[2\jot]
      \dst\frac{f(x-)+f(x+)}{2} & \mbox{if $0<x<L$ and $f$ is discontinuous at $x$ } \\[2\jot]
      0                         & \mbox{if $x=L$.}
    \end{array}\right.
  \]
\end{theorem}

\setcounter{theorem}{3}
\begin{theorem}
  \label{thmtype:11.3.4} If $f$ is piecewise smooth on $[0,L]$, then the mixed Fourier sine series
  \[
    S_{M}(x)=\sum_{n=1}^{\infty}d_{n}\sin\frac{(2n-1)\pi x}{2L}
  \]
  of $f$ on $[0,L]$, with
  \[
    d_{n}=\frac{2}{L}\int_{0}^{L}f(x)\sin\frac{(2n-1)\pi x}{2L}\,dx,
  \]
  converges for all $x$ in $[0,L]$; moreover,
  \[
    S_{M}(x)= \left\{
    \begin{array}{cl}
      0                         & \mbox{if $x=0$}                                    \\
      f(x)                      & \mbox{if $0<x<L$ and $f$ is continuous at $x$ }    \\[2\jot]
      \dst\frac{f(x-)+f(x+)}{2} & \mbox{if $0<x<L$ and $f$ is discontinuous at $x$ } \\[2\jot]
      f(L-)                     & \mbox{if $x=L$.}
    \end{array}\right.
  \]
\end{theorem}

\setcounter{theorem}{4}
\setcounter{equation}{2}
\begin{theorem}
  \label{thmtype:11.3.5} \mbox{}\\

\part{a} If $f'(0)=f'(L)=0$, $f''$ is continuous, and $f'''$ is piecewise continuous on $[0,L]$,
then
  \begin{equation}
    \label{eq:11.3.3}f(x)=a_{0}+\sum_{n=1}^{\infty}a_{n}\cos\frac{n\pi x}{L}, \quad0\le x\le L,
  \end{equation}
  with
  \begin{equation}
    \label{eq:11.3.4}a_{0}=\frac{1}{L}\int_{0}^{L}f(x)\,dx\mbox{\quad and \quad}a_{n}=\frac{2L^{2}}{n^{3}\pi^{3}}\int_{0}^{L}f'''(x)\sin\frac{n\pi x}{L}\,dx, \quad n\ge1.
  \end{equation}

  Now suppose $f'$ is continuous and $f''$ is piecewise continuous on $[0,L]$.

\part{b} If $f(0)=f(L)=0$, then
  \[
    f(x)=\sum_{n=1}^{\infty}b_{n}\sin\frac{n\pi x}{L}, \quad0\le x\le L,
  \]
  with
  \begin{equation}
    \label{eq:11.3.5}b_{n}=-\frac{2L}{n^{2}\pi^{2}}\int_{0}^{L}f''(x)\sin\frac{n\pi x}{L}\,dx.
  \end{equation}

\part{c} If $f'(0)=f(L)=0$, then
  \[
    f(x)= \sum_{n=1}^{\infty}c_{n}\cos\frac{(2n-1)\pi x}{2L}, \quad0\le x\le L,
  \]
  with
  \begin{equation}
    \label{eq:11.3.6}c_{n}=-\frac{8L}{(2n-1)^{2}\pi^{2}}\int_{0}^{L}f''(x)\cos\frac{(2n-1)\pi x}{2L}\,dx.
  \end{equation}

\part{d} If $f(0)=f'(L)=0$, then
  \[
    f(x)= \sum_{n=1}^{\infty}d_{n}\sin\frac{(2n-1)\pi x}{2L}, \quad0\le x\le L,
  \]
  with
  \begin{equation}
    \label{eq:11.3.7}d_{n}=-\frac{8L}{(2n-1)^{2}\pi^{2}}\int_{0}^{L}f''(x)\sin\frac{(2n-1)\pi x}{2L}\,dx.
  \end{equation}
\end{theorem}

\section*{Section 12.1}
\phantomsection
\addcontentsline{toc}{chapter}{Chapter 12}
\setcounter{chapter}{12}
\setcounter{section}{1}

\setcounter{theorem}{1}
\begin{theorem}
  \label{thmtype:12.1.2} A convergent infinite series
  \[
    W(z)=\sum_{n=1}^{\infty}w_{n}(z)
  \]
  can be differentiated term by term on a closed interval $[z_{1},z_{2}]$ to obtain
  \[
    W'(z)=\sum_{n=1}^{\infty}w_{n}'(z)
  \]
  (where the derivatives at $z=z_{1}$ and $z=z_{2}$ are one-sided) provided that $w_{n}'$ is
  continuous on $[z_{1},z_{2}]$ and
  \[
    |w_{n}'(z)|\le M_{n},\quad z_{1}\le z\le z_{2},\quad n=1,2,3,\dots,
  \]
  where $M_{1}$, $M_{2}$, \dots, $M_{n}$, \dots, are constants such that the series
  $\sum_{n=1}^{\infty}M_{n}$ converges.
\end{theorem}

\section*{Section 12.2}
\setcounter{chapter}{12}
\setcounter{section}{2}

\setcounter{theorem}{1}
\setcounter{equation}{8}
\begin{theorem}
  \label{thmtype:12.2.2} If $f$ and $g$ are piecewise smooth on $[0,L]$, then the series in
  \textup{(12.2.8)} converges for all $(x,t)$ and can be written as
  \begin{equation}
    \label{eq:12.2.9}u(x,t)=\frac{1}{2}[S_{f}(x+at)+S_{f}(x-at)]+ \frac{1}{2a}\int_{x-at}^{x+at}S_{g}(\tau)\,d\tau.
  \end{equation}
\end{theorem}

\setcounter{theorem}{2}
\setcounter{equation}{12}
\begin{theorem}
  \label{thmtype:12.2.3} Suppose $h$ is differentiable on $[0,L]$; that is, $h'(x)$ exists for
  $0<x<L$, and the one-sided derivatives
  \[
    h'_{+}(0)=\lim_{x\to0+}\frac{h(x)-h(0)}{x}\mbox{\quad and \quad}h'_{-}(L)=\lim_{x\to L-}\frac{h(x)-h(L)}{x-L}
  \]
  both exist.

\part{a} Let $p$ be the odd periodic extension of $h$ to $(-\infty,\infty)$; that is,
  \[
    p(x)= \left\{
    \begin{array}{cl}
      h(x),   & \quad0\le x\le L, \\
      -h(-x), & -L<x<0,
    \end{array}\right. \mbox{\  and\  }p(x+2L)=p(x),\quad-\infty<x<\infty.
  \]
  Then $p$ is differentiable on $(-\infty,\infty)$ if and only if
  \begin{equation}
    \label{eq:12.2.13}h(0)=h(L)=0.
  \end{equation}

\part{b} Let $q$ be the even periodic extension of $h$ to $(-\infty,\infty)$; that is,
  \[
    q(x)= \left\{
    \begin{array}{cl}
      h(x),  & \quad0\le x\le L, \\
      h(-x), & -L<x<0,
    \end{array}\right. \mbox{\  and\  }q(x+2L)=q(x),\quad-\infty<x<\infty.
  \]
  Then $q$ is differentiable on $(-\infty,\infty)$ if and only if
  \begin{equation}
    \label{eq:12.2.14}h'_{+}(0)=h'_{-}(L)=0.
  \end{equation}
\end{theorem}

\setcounter{theorem}{3}
\setcounter{equation}{14}
\begin{theorem}
  \label{thmtype:12.2.4} The formal solution of \textup{(12.2.1)} is an actual solution if $g$ is
  differentiable on $[0,L]$ and
  \begin{equation}
    \label{eq:12.2.15}g(0)=g(L)=0,
  \end{equation}
  while $f$ is twice differentiable on $[0,L]$ and
  \begin{equation}
    \label{eq:12.2.16}f(0)=f(L)=0
  \end{equation}
  and
  \begin{equation}
    \label{eq:12.2.17}f''_{+}(0)=f''_{-}(L)=0.
  \end{equation}
\end{theorem}

\section*{Section 13.1}
\phantomsection
\addcontentsline{toc}{chapter}{Chapter 13}
\setcounter{chapter}{13}
\setcounter{section}{1}

\setcounter{theorem}{0}
\begin{theorem}
  \label{thmtype:13.1.1} If $z_{1}$ and $z_{2}$ are solutions of $Ly=0$ such that either
  $B_{1}(z_{1})=B_{1}(z_{2})=0$ or $B_{2}(z_{1})=B_{2}(z_{2})=0$, then $\{z_{1},z_{2}\}$ is linearly
  dependent. Equivalently, if\, $\{z_{1},z_{2}\}$ is linearly independent, then
  \[
    B_{1}^{2}(z_{1})+B_{1}^{2}(z_{2})\ne0\text{\; and\; \; }B_{2}^{2}(z_{1})+B_{2}^{2}(z_{2})\ne0.
  \]
\end{theorem}

\setcounter{theorem}{1}
\setcounter{equation}{7}
\begin{theorem}
  \label{thmtype:13.1.2} The following statements are equivalent; that is, they are either all true
  or all false.
  \begin{alist}
    \item % (a)
    There's a fundamental set $\{z_{1},z_{2}\}$ of solutions of $Ly=0$ such that
    \begin{equation}
      \label{eq:13.1.8}B_{1}(z_{1})B_{2}(z_{2})-B_{1}(z_{2})B_{2}(z_{1})\ne0.
    \end{equation}

    \item % (b)
    If $\{y_{1},y_{2}\}$ is a fundamental set of solutions of $Ly=0$ then
    \begin{equation}
      \label{eq:13.1.9}B_{1}(y_{1})B_{2}(y_{2})-B_{1}(y_{2})B_{2}(y_{1})\ne0.
    \end{equation}

    \item % (c)
    For each continuous $F$ and pair of constants $(k_{1},k_{2})$, the boundary value problem
    \[
      Ly=F,\quad B_{1}(y)=k_{1},\quad B_{2}(y)=k_{2}
    \]
    has a unique solution.

    \item % (d)
    The homogeneous boundary value problem
    \begin{equation}
      \label{eq:13.1.10}Ly=0,\quad B_{1}(y)=0,\quad B_{2}(y)=0
    \end{equation}
    has only the trivial solution $y=0$.

    \item % (e)
    The homogeneous equation $Ly=0$ has linearly independent solutions $z_{1}$ and $z_{2}$ such that
    $B_{1}(z_{1})=0$ and $B_{2}(z_{2})=0$.
  \end{alist}
\end{theorem}

\setcounter{theorem}{2}
\setcounter{equation}{14}
\begin{theorem}
  \label{thmtype:13.1.3} Suppose the homogeneous boundary value problem
  \begin{equation}
    \label{eq:13.1.15}Ly=0,\quad B_{1}(y)=0,\quad B_{2}(y)=0
  \end{equation}
  has only the trivial solution. Let $y_{1}$ and $y_{2}$ be linearly independent solutions of $Ly=0$
  such that $B_{1}(y_{1})=0$ and $B_{2}(y_{2})=0$, and let
  \[
    W=y_{1}y_{2}'-y_{1}'y_{2}.
  \]
  Then the unique solution of
  \begin{equation}
    \label{eq:13.1.16}Ly=F,\quad B_{1}(y)=0,\quad B_{2}(y)=0
  \end{equation}
  is
  \begin{equation}
    \label{eq:13.1.17}y(x)=y_{1}(x)\int_{x}^{b}\frac{F(t)y_{2}(t)}{P_{0}(t)W(t)}\,dt+ y_{2}(x)\int_{a}^{x}\frac{F(t)y_{1}(t)}{P_{0}(t)W(t)}\,dt.
  \end{equation}
\end{theorem}

\setcounter{theorem}{3}
\setcounter{equation}{21}
\begin{theorem}
  \label{thmtype:13.1.4} Suppose the homogeneous boundary value problem
  \begin{equation}
    \label{eq:13.1.22}Ly=0,\quad B_{1}(y)=0,\quad B_{2}(y)=0
  \end{equation}
  has a nontrivial solution $y_{1}$, and let $y_{2}$ be any solution of $Ly=0$ that isn't a constant
  multiple of $y_{1}$. Let $W=y_{1}y_{2}'-y_{1}'y_{2}$. If
  \begin{equation}
    \label{eq:13.1.23}\int_{a}^{b}\frac{F(t)y_{1}(t)}{P_{0}(t)W(t)}\,dt=0,
  \end{equation}
  then the nonhomogeneous boundary value problem
  \begin{equation}
    \label{eq:13.1.24}Ly=F,\quad B_{1}(y)=0,\quad B_{2}(y)=0
  \end{equation}
  has infinitely many solutions, all of the form $y=y_{p}+c_{1}y_{1}$, where
  \[
    y_{p}=y_{1}(x) \int_{x}^{b}\frac{F(t)y_{2}(t)}{P_{0}(t)W(t)}\,dt+ y_{2}(x) \int_{a}^{x}\frac{F(t)y_{1}(t)}{P_{0}(t)W(t)}\,dt
  \]
  and $c_{1}$ is a constant. If
  \[
    \int_{a}^{b}\frac{F(t)y_{1}(t)}{P_{0}(t)W(t)}\,dt\ne0,
  \]
  then \eqref{eq:13.1.24} has no solution.
\end{theorem}

\section*{Section 13.2}
\setcounter{chapter}{13}
\setcounter{section}{2}

\setcounter{theorem}{0}
\setcounter{equation}{4}
\begin{theorem}
  \label{thmtype:13.2.1} If $P_{0}$, $P_{1}$, $P_{2}$, and $R$ are continuous and $P_{0}$ and $R$
  are positive on a closed interval $[a,b]$, then the equation
  \begin{equation}
    \label{eq:13.2.5}P_{0}(x)y''+P_{1}(x)y'+P_{2}(x)y+\lambda R(x)y=0
  \end{equation}
  can be rewritten as
  \begin{equation}
    \label{eq:13.2.6}(p(x)y')'+q(x)y+\lambda r(x)y=0,
  \end{equation}
  where $p$, $p'$, $q$ and $r$ are continuous and $p$ and $r$ are positive on $[a,b]$.
\end{theorem}

\setcounter{theorem}{1}
\setcounter{equation}{20}
\begin{theorem}
  \label{thmtype:13.2.2} If
  \[
    Ly=(p(x)y')'+q(x)y
  \]
  and $u$ and $v$ are twice continuously differentiable functions on $[a,b]$ that satisfy the boundary conditions
  $B_{1}(y)=0$ and $B_{2}(y)=0$, then
  \begin{equation}
    \label{eq:13.2.21}\int_{a}^{b}[u(x)Lv(x)-v(x)Lu(x)]\,dx=0.
  \end{equation}
\end{theorem}

\setcounter{theorem}{2}
\begin{theorem}
  \label{thmtype:13.2.3} If $\lambda=p+qi$ with $q\ne0$ then the boundary value problem
  \[
    Ly+\lambda r(x)y=0,\quad B_{1}(y)=0,\quad B_{2}(y)=0
  \]
  has only the trivial solution.
\end{theorem}

\setcounter{theorem}{3}
\setcounter{equation}{24}
\begin{theorem}
  \label{thmtype:13.2.4} If $\lambda_{1}$ and $\lambda_{2}$ are distinct eigenvalues of the
  Sturm--Liouville problem
  \begin{equation}
    \label{eq:13.2.25}Ly+\lambda r(x)y=0,\quad B_{1}(y)=0,\quad B_{2}(y)=0
  \end{equation}
  with associated eigenfunctions $u$ and $v$ respectively, then
  \begin{equation}
    \label{eq:13.2.26}\int_{a}^{b}r(x)u(x)v(x)\,dx=0.
  \end{equation}
\end{theorem}

\setcounter{theorem}{4}
\begin{theorem}
  \label{thmtype:13.2.5} If $u\not\equiv0$ and $v$ both satisfy
  \[
    Ly+\lambda r(x)y=0,\quad B_{1}(y)=0,\quad B_{2}(y)=0,
  \]
  then $v=cu$ for some constant $c$.
\end{theorem}

\setcounter{theorem}{5}
\setcounter{equation}{26}
\begin{theorem}
  \label{thmtype:13.2.6} The set of all eigenvalues of the Sturm--Liouville problem
  \[
    Ly+\lambda r(x)y=0,\quad B_{1}(y)=0,\quad B_{2}(y)=0
  \]
  can be ordered as
  \[
    \lambda_{1}<\lambda_{2}<\cdots<\lambda_{n}<\cdots,
  \]
  and
  \[
    \lim_{n\to\infty}\lambda_{n}=\infty.
  \]
  For each $n$, if $y_{n}$ is an arbitrary $\lambda_{n}$-eigenfunction, then every
  $\lambda_{n}$-eigenfunction is a constant multiple of $y_{n}$. If $m\ne n$, $y_{m}$ and $y_{n}$
  are orthogonal on $[a,b]$ with respect to $r=r(x)$; that is,
  \begin{equation}
    \label{eq:13.2.27}\int_{a}^{b}r(x)y_{m}(x)y_{n}(x)\,dx=0.
  \end{equation}
\end{theorem}

\setcounter{theorem}{6}
\begin{theorem}
  \label{thmtype:13.2.7} Let $\lambda_{1}<\lambda_{2}<\cdots<\lambda_{n}<\cdots$ be the eigenvalues
  of the Sturm--Liouville problem
  \[
    Ly+\lambda r(x)y=0,\quad B_{1}(y)=0,\quad B_{2}(y)=0,
  \]
  with associated eigenfunctions $y_{1}$, $y_{2}$, \dots, $y_{n}$, \dots. Suppose $f$ is piecewise
  smooth (Definition~11.2.3) on $[a,b]$. For each $n$, let
  \[
    c_{n}=\frac{\dst{\int_{a}^{b}}r(x)f(x)y_{n}(x)\,dx}{\dst{\int_{a}^{b}}r(x)y_{n}^{2}(x)\,dx}.
  \]
  Then
  \[
    \frac{f(x-)+f(x+)}{2}=\sum_{n=1}^{\infty}c_{n}y_{n}(x)
  \]
  for all $x$ in the open interval $(a,b)$.
\end{theorem}

\end{document}
